\documentclass[11pt]{article}
% Shared preamble for the rigor documents (Lamport-style structured proofs).  \input this file after \documentclass.
\usepackage[T1]{fontenc}\usepackage{lmodern}\usepackage{microtype}
\usepackage[margin=1in]{geometry}
\usepackage{amsmath,amssymb,amsthm,mathtools}
\usepackage{xcolor}\usepackage{enumitem}\usepackage{booktabs}\usepackage{graphicx}\usepackage{hyperref}
\usepackage[most]{tcolorbox}
\definecolor{navy}{HTML}{17365D}\definecolor{teal}{HTML}{117A8B}\definecolor{warn}{HTML}{A33A2B}\definecolor{okgreen}{HTML}{2E7D4F}
\hypersetup{colorlinks=true,linkcolor=navy,citecolor=teal,urlcolor=teal}
\theoremstyle{plain}
\newtheorem{theorem}{Theorem}[section]\newtheorem{lemma}[theorem]{Lemma}\newtheorem{proposition}[theorem]{Proposition}\newtheorem{corollary}[theorem]{Corollary}
\theoremstyle{definition}\newtheorem{definition}[theorem]{Definition}\newtheorem{remark}[theorem]{Remark}\newtheorem{claim}[theorem]{Claim}\newtheorem{issue}[theorem]{Issue}
% ---------- Lamport structured proofs ----------
% \lstep{level}{number}{assertion}   -> indented "<level>number. assertion"
% \lproof                             -> "PROOF:" line (start of the sub-proof of the preceding step)
% \lby{justification}                 -> "BY ..." leaf justification for the preceding step
% \lqed{level}{number}                -> QED step of a sub-proof
% keywords: \lassume \lprove \lsuffices \lcase \ldefine \lpick \lnew
\newlength{\lindent}\setlength{\lindent}{1.7em}
\newcommand{\lstep}[3]{\par\smallskip\noindent\hspace*{\dimexpr\numexpr#1-1\relax\lindent\relax}{\color{navy}$\langle#1\rangle#2$.}\ #3}
\newcommand{\lproof}{\par\noindent\hspace*{\lindent}{\scshape Proof:}}
\newcommand{\lby}[1]{\par\noindent\hspace*{\lindent}{\scshape By}\ #1.}
\newcommand{\lqed}[2]{\par\smallskip\noindent\hspace*{\dimexpr\numexpr#1-1\relax\lindent\relax}{\color{navy}$\langle#1\rangle#2$.}\ {\scshape Q.E.D.}}
\newcommand{\lassume}{{\scshape Assume}\ }\newcommand{\lprove}{{\scshape Prove}\ }\newcommand{\lsuffices}{{\scshape Suffices}\ }
\newcommand{\lcase}{{\scshape Case}\ }\newcommand{\ldefine}{{\scshape Define}\ }\newcommand{\lpick}{{\scshape Pick}\ }\newcommand{\lnew}{{\scshape New}\ }
% ---------- verdict boxes ----------
\newtcolorbox{verdictok}{colback=okgreen!8,colframe=okgreen,title={\bfseries Verdict: verified},fonttitle=\small}
\newtcolorbox{verdictgap}{colback=orange!10,colframe=orange!80!black,title={\bfseries Verdict: gap / caveat},fonttitle=\small}
\newtcolorbox{verdictbreak}{colback=warn!10,colframe=warn,title={\bfseries Verdict: proof-breaking issue},fonttitle=\small}
\newtcolorbox{citedbox}[1]{colback=navy!5,colframe=navy!60,title={\bfseries Cited result: #1},fonttitle=\small,breakable}
\setlength{\parindent}{0pt}\setlength{\parskip}{4pt}\allowdisplaybreaks
\newcommand{\cA}{\mathcal A}\newcommand{\cF}{\mathcal F}\newcommand{\cM}{\mathfrak M}\newcommand{\cN}{\mathfrak N}

% extra calligraphic shortcuts (\providecommand: no clash if a document defines its own)
\providecommand{\cB}{\mathcal B}\providecommand{\cC}{\mathcal C}\providecommand{\cD}{\mathcal D}\providecommand{\cE}{\mathcal E}\providecommand{\cG}{\mathcal G}\providecommand{\cH}{\mathcal H}\providecommand{\cI}{\mathcal I}\providecommand{\cJ}{\mathcal J}\providecommand{\cK}{\mathcal K}\providecommand{\cL}{\mathcal L}\providecommand{\cO}{\mathcal O}\providecommand{\cP}{\mathcal P}\providecommand{\cQ}{\mathcal Q}\providecommand{\cR}{\mathcal R}\providecommand{\cS}{\mathcal S}\providecommand{\cT}{\mathcal T}\providecommand{\cU}{\mathcal U}\providecommand{\cV}{\mathcal V}\providecommand{\cW}{\mathcal W}\providecommand{\cX}{\mathcal X}\providecommand{\cY}{\mathcal Y}\providecommand{\cZ}{\mathcal Z}
\newcommand{\Fid}{\operatorname{F}}\newcommand{\Tr}{\operatorname{Tr}}\newcommand{\eps}{\varepsilon}\newcommand{\dd}{\,\mathrm d}

\newcommand{\Sch}{\mathcal S}
\newcommand{\Mob}{\operatorname{M\ddot ob}}
\newcommand{\bI}{\beta_I}
\newcommand{\id}{\mathrm{id}}
\newcommand{\Rp}{\mathsf R}
\newcommand{\Lp}{\mathsf L}
\newcommand{\Fr}{\cF_r}
\newcommand{\Bang}{\mathsf A}
\newcommand{\PP}{\mathcal P}
\newcommand{\Ad}{\operatorname{Ad}}
\theoremstyle{definition}\newtheorem{example}[theorem]{Example}
\newtheorem{hypothesis}[theorem]{Hypothesis}

\title{The corner calculus of sequential one-sided recovery:\\
Schwarzian point masses, the recovered state, the separation constraint,\\
and the type-III sequential bound}
\author{Track G1a (Phase 6, analytic)}
\date{2026-09-16}

\begin{document}
\maketitle

\begin{abstract}
A \emph{one-sided protocol} on a chain of adjacent intervals $A_1,\dots,A_n$ recovers the vacuum
on $I=A_1\cup\dots\cup A_n$ from a starting block by a sequence of zero-collar (rotated or
ordinary) Petz steps, each of which is a parabolic M\"obius compression fixing a junction.  We
prove: (i) the first-order field of a step is \emph{exactly} $\zeta\,w(y-y_p)\partial_y$ in the
modular frame of any interval containing the corner, $w(u)=-4\sinh^2(u/2)\mathbf 1_{u>0}$,
$\zeta=\sigma\bI(p)$ (Lemma~\ref{lem:field}); (ii) the composite point map $\Phi$ is $C^1$
piecewise-M\"obius and its distributional Schwarzian is a sum of point masses, located at the
junctions with weights $-2\kappa_k$ under a hypothesis (H) on the protocol, and in general at the
preimages of the junctions with weights scaled by the derivative of the later maps
(Theorem~\ref{thm:corner}); we exhibit two failures of (H), one of which shows that
\emph{single-interval conditioning alone does not imply} (H) --- a correction to the wording of
W1; (iii) the recovered state depends on $\Phi$ only through $\Sch(\Phi)$
(Theorem~\ref{thm:state}), it equals the vacuum if and only if $\Phi$ is M\"obius
(Theorem~\ref{thm:rigidity}), so no protocol whose last step keeps a non-empty part recovers
exactly; (iv) VWZ Protocol 1(B)
$=$ 1(A) as point maps, for every common rotation parameter, and VWZ Protocol 3 in either order
is, up to a global M\"obius map, the over-compressed single step of Theorem~A with
$\zeta_{\rm eff}=(\sigma_1+\sigma_2)\bI(p_3)$ (\S\ref{sec:vwz}); (v) with single-interval
conditioning the corner measure --- hence the recovered state --- depends only on the
\emph{starting block}, not on the order of the steps, which gives the complete tables for $n=4,5$
(\S\ref{sec:tables}); (vi) the corner strength in the frame of the whole chain is
$\zeta^{(I)}=\zeta_{\rm step}(1+z')$ (Lemma~\ref{lem:ampl}); (vii) for the left-to-right Petz
chain the corner data take the purely modular form
$\zeta_k=\sinh(\Delta_{k+1}/2)/[\sinh(\Delta_k/2)\sinh((\Delta_k+\Delta_{k+1})/2)]$, from which we
\emph{prove} the separation bound $e^{\Delta_k}\ge2/\eps$ for all-weak protocols
(Theorem~\ref{thm:sep}), sharp as $\eps\to0$ (the constant $2$ is the $\eps\to0$ limit of the exact numerical minimum);
(viii) the Bures-angle triangle inequality gives the type-III sequential bound
$\Bang_{\rm tot}\le\sum_k\Bang_k$ for any chain of von Neumann inclusions
(Theorem~\ref{thm:w6}) under Hypothesis~(U), and the chordal Bures-distance route gives the
second-order form $\Phi_{\rm tot}\le(\sum_k\sqrt{\Phi_k})^2(1+o(1))$ with no hypothesis
(Theorem~\ref{thm:w6chord}, Corollary~\ref{cor:w6-second}(2)).  Every closed-form identity is verified numerically in
\texttt{rigor/g1a\_calculus.py} and \texttt{rigor/g1a\_separation.py}.
\end{abstract}

\tableofcontents

\section{Setting, conventions, and the class of protocols}\label{sec:setting}

Every convention is fixed here once and cited afterwards.  They are the conventions of
\texttt{rigor/phase6\_brief.md} (\S``Setting and conventions''), which in turn follow
paper~1 (\texttt{paper1/sec\_setting.tex}, \texttt{paper1/sec\_petz.tex}) and Note~4
(\texttt{continuum\_petz\_free\_fermion.tex}).

\paragraph{(C1) Geometry.}
Fix $n\ge2$ and points $p_1<p_2<\dots<p_{n+1}$ on the chiral light ray $\mathbb R$.  Put
\[
 A_k:=(p_k,p_{k+1}),\quad l_k:=p_{k+1}-p_k>0,\quad
 L_k:=p_{k+1}-p_1=l_1+\dots+l_k,\quad T:=L_n,\quad I:=(p_1,p_{n+1}).
\]
The $p_k$ with $2\le k\le n$ are the \emph{junctions}.  A \emph{block} is a union
$A_i\cup A_{i+1}\cup\dots\cup A_j$ of consecutive chain intervals, viewed as the open interval
$(p_i,p_{j+1})$ (the interior junctions are not removed; the local algebras of the interval and
of the union agree by strong additivity, \texttt{paper1/sec\_setting.tex}, l.~40--42).

\paragraph{(C2) Net, vacuum, one-particle space.}
$\Fr$ is the net of $r$ chiral free fermions on $\mathbb R$ (Longo--Xu \cite{LX} \S3.1; $c=r$),
$\omega$ its vacuum state, $H:=L^2(\mathbb R;\mathbb C^r)$, $P$ the Hardy projection,
$H_J:=L^2(J;\mathbb C^r)$ and $Q_J:=E_JPE_J$ the vacuum symbol of an interval $J$.  For a
bounded operator $V$ on $H$ preserving $H_J$ we write $\beta_V$ for the Bogoliubov
$*$-homomorphism determined by $\beta_V(\psi(f))=\psi(Vf)$; $\beta_{V_1}\beta_{V_2}=\beta_{V_1V_2}$.
For an orientation-preserving $C^1$ diffeomorphism $k:J\to k(J)\subseteq\mathbb R$ we write
\begin{equation}\label{eq:pushforward}
 (V_kf)(y):=\bigl((k^{-1})'(y)\bigr)^{1/2}f\bigl(k^{-1}(y)\bigr)\ \ (y\in k(J)),\qquad
 V_kf:=0 \text{ off } k(J),
\end{equation}
the \emph{half-density push-forward}; $V_k:H_J\to H_{k(J)}$ is an isometry and
\begin{equation}\label{eq:functorial}
 V_{k_1}V_{k_2}=V_{k_1\circ k_2}
\end{equation}
whenever the composition is defined (\S\ref{sec:state}, Lemma~\ref{lem:comp-W}).  All maps
are orientation preserving, so no signs or phases occur.

\paragraph{(C3) The one-sided step.}
Let $J$ be a block and $A_{\rm new}$ a chain interval adjacent to $J$ and disjoint from it.  A
\emph{one-sided step} adjoining $A_{\rm new}$ to $J$ is specified by a \emph{conditioning region}
$A_B\subseteq J$, which is a block of $J$ containing the endpoint of $J$ shared with
$A_{\rm new}$, together with a parameter $s$.  Write $l_B:=|A_B|$, $l_{\rm new}:=|A_{\rm new}|$,
$\lambda:=l_{\rm new}/l_B$ and
\begin{equation}\label{eq:corner-def}
 \text{\emph{corner}}\quad p:=\text{the endpoint of $A_B$ \emph{not} shared with $A_{\rm new}$},
 \qquad
 M:=A_B\cup A_{\rm new}\ \ (\text{the \emph{moving part}}),
\end{equation}
an interval of length $l_B+l_{\rm new}$ with $p$ as one endpoint.  If $A_B\subsetneq J$ then $p$
is the junction of $A_B$ with the rest of $J$; the degenerate case $A_B=J$ (empty kept part) is
allowed, and then $p$ is the far endpoint of $J$.  We say a step \emph{keeps a non-empty part} if
$A_B\subsetneq J$; this is Note~4's hypothesis $a>0$ (Box~[G1]).  With
\begin{equation}\label{eq:sigma}
 \sigma:=\frac s{l_B+l_{\rm new}},\qquad s\in[\lambda,2\lambda],
\end{equation}
the \emph{step map} $k:J\cup A_{\rm new}\to J$ is
\begin{equation}\label{eq:step-map}
 k=\begin{cases}
 \Rp_{p,\sigma}:\ x\mapsto x\ (x\le p),\quad x\mapsto p+\dfrac{x-p}{1+\sigma(x-p)}\ (x>p),
 &\text{$A_{\rm new}$ to the right of $p$},\\[2.2ex]
 \Lp_{p,\sigma}:\ x\mapsto x\ (x\ge p),\quad x\mapsto p+\dfrac{x-p}{1-\sigma(x-p)}\ (x<p),
 &\text{$A_{\rm new}$ to the left of $p$},
 \end{cases}
\end{equation}
restricted to $J\cup A_{\rm new}$.  It is the identity on $J\setminus A_B$ and compresses $M$
towards $p$.  $s=2\lambda$ is the ordinary Petz map (VWZ's $\lambda=0$) and
$s=s_t=\lambda(1+e^{-2\pi t})$ the rotated ones (\texttt{paper1/sec\_petz.tex}
Def.~$\zeta$, Box~[G2] below; Note~4 Thm.~7.1).  The condition $s\ge\lambda$ is exactly
$k(M)\subseteq A_B$ (Lemma~\ref{lem:admissible}), which is what Note~4 Thm.~7.1 requires.
Two conventions that matter and were left tacit in the first version (REF-P6-3 \S7, \S11).
(i) \emph{Each step is the (rotated) Petz map of the \emph{vacuum} on its own inclusion}
$\Fr(A_B)\subset\Fr(A_B\cup A_{\rm new})$, applied to whatever state the previous steps have
produced --- not the Petz map of that intermediate state.  This is also what VWZ do (their
$P_{R\to RS}$ is built from the vacuum reduced density matrices and then applied to
$\sigma^{(1)}$, Box~[G6]); the adaptive variant is a different channel and is not treated here.
(ii) $s\in[\lambda,2\lambda]$ corresponds to $t\in[0,\infty]$.  Nothing below needs the upper
end: Note~4 Thm.~7.1 needs only $s\ge\lambda$ and the corner calculus only $\sigma>0$, so the
$t$-average of \S\ref{sec:twirl} over all of $\mathbb R$ (where $s_t>2\lambda$ and
$\theta_t=\frac12(1+e^{-2\pi t})>1$ for $t<0$) is admissible in that wider sense; statements
that quote $\theta_t\in(\frac12,1]$ are for $t\ge0$.

\paragraph{(C4) Protocols and the composite point map.}
A \emph{one-sided protocol} $\PP$ consists of a starting block $J_0$ and steps
$m=1,\dots,N$, the $m$-th adjoining a chain interval $A^{(m)}_{\rm new}$ to $J_{m-1}$ with
conditioning region $A^{(m)}_B$, corner $p^{(m)}$, moving part $M_m$, parameter $\sigma_m>0$ and
step map $k_m:J_m\to J_{m-1}$, where $J_m:=J_{m-1}\cup A^{(m)}_{\rm new}$; the protocol is
\emph{complete} if $J_N=I$.  The \emph{composite point map} is
\begin{equation}\label{eq:Phi}
 \Phi:=k_1\circ k_2\circ\dots\circ k_N:\ J_N\longrightarrow J_0
 \qquad(\text{the \emph{first} step outermost}),
\end{equation}
and the \emph{corner measure} of $\PP$ is the positive measure
$\kappa_\PP:=\sum_m\sigma_m\delta_{p^{(m)}}$; we write $\kappa_k$ for the mass at the junction
$p_k$.  Order \eqref{eq:Phi} is forced by (C2): the Heisenberg-picture step map is
$\beta_m:=\beta_{V_{k_m}}:\Fr(J_m)\to\Fr(J_{m-1})$, so the composite recovery map is
$\beta_1\circ\beta_2\circ\dots\circ\beta_N=\beta_{V_\Phi}$ by \eqref{eq:functorial}, and the
recovered state on $\Fr(J_N)$ is $\omega\circ\beta_{V_\Phi}$.

\paragraph{(C5) Modular frame.}
For an interval $I=(p_1,p_{n+1})$ put
\begin{equation}\label{eq:modframe}
 y=y_I(x):=\log\frac{x-p_1}{p_{n+1}-x},\qquad
 \bI(x):=\frac{\dd x}{\dd y}=\frac{(x-p_1)(p_{n+1}-x)}{T} .
\end{equation}
$y_I$ maps $I$ onto $\mathbb R$ and is the modular coordinate of $I$ (the dilation generator of
$I$ is $\partial_y$; BGL \cite{BGL} Thm.~2.3(ii), Box~[G3]).  For a corner $(p,\kappa)$
we set $\zeta:=\kappa\,\bI(p)$ and $y_p:=y_I(p)$.

\paragraph{(C6) Schwarzian conventions.}
For $\Phi$ with $\Phi'>0$ and $v:=(\log\Phi')'=\Phi''/\Phi'$ of locally bounded variation we
define the Schwarzian derivative as the distribution
\begin{equation}\label{eq:schwarz}
 \Sch(\Phi):=v'-\tfrac12v^2
 \qquad\Bigl(=\frac{\Phi'''}{\Phi'}-\frac32\Bigl(\frac{\Phi''}{\Phi'}\Bigr)^2
 \text{ where $\Phi$ is $C^3$}\Bigr);
\end{equation}
$v'$ is a measure and $v^2\in L^\infty_{\rm loc}$, so \eqref{eq:schwarz} is well defined
(\S\ref{sec:schwarz}).  $\Mob$ denotes the group of orientation-preserving real M\"obius maps
$x\mapsto(ax+b)/(cx+d)$, $ad-bc>0$.

\paragraph{(C7) Fidelity conventions.}
$P_{\rm tr}(\varphi,\psi)$ is Uhlmann's transition probability (Box~[G4]),
$F:=P_{\rm tr}^{1/2}$ the fidelity, $\Bang:=\arccos F\in[0,\pi/2]$ the \emph{Bures angle}, and
$\Phi_{\rm fid}:=-\log F$ the quantity called $\Phi$ in paper~1.  We write $\Phi_{\rm fid}$ with
its subscript throughout, because $\Phi$ denotes the point map \eqref{eq:Phi}.  $\Phi_A(\zeta)$
is paper~1's single-step function: at finite $\zeta$, $\Phi_A(\zeta)$ is a \emph{definition}
--- $-\log F$ of the vacuum against the single geometric compression with invariant $\zeta$,
well posed by Corollary~\ref{cor:phiA} --- and the \emph{theorem} attached to it is the
$\zeta\to0$ law $\Phi_A(\zeta)=f_2\zeta^2+o(\zeta^2)$, $f_2=c/(12\pi^2)$ (card
\texttt{thm-a-quadratic-law}); paper~1's tables give it numerically at finite $\zeta$.

\section{Cited results}\label{sec:cited}

Each box states what is used, where it comes from, and how this document uses it, following
\texttt{rigor/README\_AGENTS.md}.

\begin{citedbox}{[G1] Note~4 Thm.~7.1 $=$ paper~1 Thm.~2.6 (\texttt{thm:normal}): the zero-collar
step map is a normal $*$-isomorphism onto its range}
\emph{Transcribed verbatim} (\texttt{continuum\_petz\_free\_fermion.tex}, Thm.~7.1, p.~12 of the
compiled PDF; identical statement in \texttt{paper1/sec\_petz.tex} Thm.~``normal'', p.~6):
``For every $s\ge\lambda$, the $C^*$-isomorphism $\beta_s$ in (beta-Cstar) extends
uniquely to a normal, parity-preserving $*$-isomorphism onto its range,
$\overline\beta_s:\cF_r(I)\to\cF_r(J_s)\subset\cF_r(AB)$.  For $s=s_t$, its restriction to every
positive-collar algebra is the Faulkner--Hollands--Swingle--Wang rotated Petz map,
$\overline\beta_{s_t}|_{\cM_\eps}=\alpha_{\eps,t}$.''
Here $A=(-a,0)$, $B=(0,b)$, $C=(b,b+c)$, $I=ABC$, $L=b+c$, $\lambda=c/b$, $J_s=(-a,L/(1+s))$, and
$\beta_s$ is the geometric compression $h_s(x)=x/(1+sx/L)$ on $D=BC$, $\id$ on $A$ (Note~4
eqs.~(h-s), (k-s); \texttt{paper1/sec\_setting.tex} l.~62--66).  The proof uses the exact
trace-norm identity $\lVert Q_I-\widetilde Q_s\rVert_1=\frac r{2\pi}\log(1+\frac{as}{a+L})$
(Note~4 Thm.~6.1), the Powers--St\o rmer inequality and the Powers--St\o rmer--Araki
quasi-equivalence criterion.

\emph{Use here.}  Each individual step $\beta_m$ of a protocol is of this form, with
$A=J_{m-1}\setminus A^{(m)}_B$, $B=A^{(m)}_B$, $C=A^{(m)}_{\rm new}$, $I=J_m$; the theorem gives
normality of $\beta_m$, which we compose in Theorem~\ref{thm:state}(1).  \emph{It is not applied
to $\Phi$ itself}: $\Phi$ is not the zero-collar map of a single inclusion.

\emph{Verdict on the usage.}  Hypotheses satisfied: $s_m\ge\lambda_m$ by (C3), the three intervals
are adjacent with the corner as the $A|B$ junction, and the step acts as the identity on
$J_{m-1}\setminus A^{(m)}_B$ exactly as $\beta_s$ acts as the identity on $A$.  One caveat is
recorded in Remark~\ref{rem:caveat}.

\smallskip\noindent\shot[\textwidth]{G1a_N4_thm71.png}
\end{citedbox}

\begin{citedbox}{[G2] paper~1 Def.~2.5 (\texttt{def:zeta}): the invariant $\zeta$}
\emph{Transcribed verbatim} (\texttt{paper1/sec\_petz.tex} l.~105--108, p.~6 of
\texttt{paper1/main.pdf}): ``\textbf{Definition.} $\zeta:=\frac{as}{a+L}$.  For the ordinary Petz
map $s=2c/b$ and $\zeta=2z$; for the rotated map $s_t$ one has $\zeta_t=z(1+e^{-2\pi t})$.''
Here $z=ac/(b(a+b+c))$ is the cross ratio of \texttt{paper1/sec\_setting.tex} eq.~(eta-z).

\emph{Use here.}  Lemma~\ref{lem:field} identifies $\zeta$ with $\sigma\bI(p)$ for the frame
$I=(-a,L)$, $p=0$, $\sigma=s/L$, and Definition~\ref{def:strength} extends it to an arbitrary
frame.  Theorem~\ref{thm:prot3} produces a $\zeta_{\rm eff}$ of exactly this form.

\smallskip\noindent\shot[\textwidth]{G1a_P1_defzeta_thmnormal.png}
\end{citedbox}

\begin{citedbox}{[G3] BGL 1993 Thm.~2.3(ii): geometric modular group; M\"obius invariance of the
vacuum}
\emph{Transcribed verbatim} (Brunetti--Guido--Longo \cite{BGL}, Thm.~2.3, arXiv
funct-an/9302008; screenshot \texttt{shots/R1\_BGL\_thm23.png} in \texttt{cited\_R1.tex} Box~C10):
``\textbf{2.3 Theorem.} Let $A$ be a conformally covariant pre-cosheaf on $M$, $\widetilde A$ the
corresponding pre-cosheaf on $\widetilde M$.  Then: \dots\ (ii) If $O$ is a region in
$\widetilde{\mathcal K}$, $\Delta_O$ the modular operator for $(R(O),\Omega)$ and $U_O(t)$ the
unitary associated with $\Lambda^O_t$, then $\Delta^{it}_O=U_O(t)$.''  Conformal covariance
includes $U(g)\Omega=\Omega$ for $g$ in the M\"obius group, hence
$\omega\circ\Ad U(g)=\omega$.

\emph{Use here.}  Two places.  (a) (C5): the modular flow of $I$ is the dilation
$\partial_y$ in the coordinate \eqref{eq:modframe}, which is what makes $\zeta=\kappa\bI(p)$ the
frame-covariant strength.  (b) Theorem~\ref{thm:state}(3): $\omega\circ\beta_g=\omega$ for
$g\in\Mob$.  For the free fermion we do not need to import (b): it is proved from scratch at the
one-particle level in Lemma~\ref{lem:moebius-inv}, because the M\"obius group preserves the Hardy
space.  BGL is cited for the general Möbius-covariant net, where Lemma~\ref{lem:moebius-inv} is
replaced by the covariance axiom.

\emph{Verdict.}  Used only through the two consequences just listed, both of which are also
available directly for $\Fr$ (Lemma~\ref{lem:moebius-inv}); no hypothesis of BGL is in question.

\smallskip\noindent\shot[\textwidth]{R1_BGL_thm23.png}
\end{citedbox}

\begin{citedbox}{[G4] Alberti--Uhlmann 2002: transition probability and its variational
formula}
\emph{Transcribed verbatim} (\cite{AU02}, Def.~2, p.~4, and eq.~(1-4)):
``$P_M(\nu,\varrho)=\sup_{\{\pi,K\},\phi\in S_{\pi,M}(\nu),\psi\in S_{\pi,M}(\varrho)}
|\langle\psi,\phi\rangle|^2$'', where $S_{\pi,M}(\mu)=\{\xi\in K:\mu(\cdot)=\langle\pi(\cdot)\xi,
\xi\rangle\}$; and ``$|f(1)|^2\le P_M(\nu,\varrho)\le\nu(a)\varrho(a^{-1})$'' for every invertible
$a\in M_+$.  Two exact variational formulas, with their correct locators (REF-P6-3 \S9):
\textsc{Theorem 2}(1), \textbf{p.~6}, the \emph{product} form,
``$\sqrt{P_M(\nu,\varrho)}=\inf_{x>0}\sqrt{\nu(x)\varrho(x^{-1})}$'', equivalently
$P_M(\nu,\varrho)=\inf_{a>0}\nu(a)\varrho(a^{-1})$; and \textsc{Corollary 2}(5), \textbf{p.~8},
the \emph{normalised} form,
``$\sqrt{P_M(\nu,\varrho)}=\inf_{x>0}\frac12\{\nu(x)+\varrho(x^{-1})\}$''.  (An earlier version
of this box attributed the normalised form to Thm.~2(1) and reproduced p.~11, which is the
\emph{proof} of Cor.~2(5); the screenshots below are the statements.)

\emph{Use here.}  Monotonicity of $F$ under normal unital Schwarz maps
(Lemma~\ref{lem:monotone}), concavity of $F$ in one argument (\S\ref{sec:twirl}), and, together
with Hypothesis~(U) of \S\ref{sec:w6}, the Bures-angle triangle inequality
(Theorem~\ref{thm:w6}).

\smallskip\noindent\shot[\textwidth]{G1a_AU02_thm2_p6.png}

\smallskip\noindent\shot[\textwidth]{G1a_AU02_cor2_p8.png}
\end{citedbox}

\begin{citedbox}{[G5] FHSW 2020: rotated Petz maps for half-sided modular
inclusions}
\emph{Faithful paraphrase} (Faulkner--Hollands--Swingle--Wang \cite{FHSW}, \S4.2; screenshots
\texttt{shots/FHSW\_hsm.png}, \texttt{shots/FHSW\_petz.png} in \texttt{cited\_R2.tex}): for a
half-sided modular inclusion $\cN\subset\cM$ with common cyclic separating vector, the rotated
Petz map $\alpha_t$ of the inclusion acts geometrically as a Borchers translation
proportional to $1+e^{-2\pi t}$ (made precise below); the twirled Petz map is the average $\int\dd t\,p(t)\alpha_t$ with the normalised probability
density of FHSW eq.~(26), p.~6,
\[
 p(t)=\frac\pi{\cosh(2\pi t)+1},\qquad \int_{\mathbb R}p(t)\dd t=1
\]
(transcribed verbatim: ``$p(t)$ is the normalized probability density defined by
$p(t)=\frac\pi{\cosh(2\pi t)+1}$'').  In the geometric form FHSW eq.~(167), p.~31, reads
$\alpha^t_\eta(x)=U\bigl(a(1+e^{-2\pi t})\bigr)^*x\,U\bigl(a(1+e^{-2\pi t})\bigr)$ for
$\cA_a\subset\cA$, i.e.\ the translation parameter is $a(1+e^{-2\pi t})$, which in the
normalisation where the ordinary Petz map ($t=0$) translates by $2a$ is the family $s_t$ of
(C3).

\emph{Use here.}  It is the source of the family $s_t=\lambda(1+e^{-2\pi t})$ in (C3) (via Note~4,
which proves that for $\Fr(B)\subset\Fr(BC)$ these maps are the compressions \eqref{eq:step-map}),
of the twirled channel of \S\ref{sec:twirl}, and of the general type-III recovery setting of
\S\ref{sec:w6}.

\smallskip\noindent\shot[\textwidth]{G1a_FHSW_pt_eq26.png}

\smallskip\noindent\shot[\textwidth]{FHSW_hsm.png}
\end{citedbox}

\begin{citedbox}{[G6] VWZ 2023 \S5: the four protocols and the numerical claim}
\emph{Transcribed verbatim} (Vardhan--Wei--Zou \cite{VWZ}, \S5, pp.~43--48, eqs.~(5.9)--(5.11)):
``In Protocol 1(B), we act with $P_{B\to BC}$ in the first step, forming the output state
$\sigma^{(1)}_{ABC}$.  In the next step, we act with $P_{BC\to BCD}$ on $\sigma^{(1)}_{ABC}$.
While this procedure naively appears different from Protocol 1(A), we can immediately see from the
definition of $P_{R\to RS}$ that it gives rise to exactly the same final state on $ABCD$'';
``In Protocol 2, we start with the state $\rho_{AB}$, and the first step is the same as in
Protocol 1(B).  In the next step, we act with the map $P_{C\to CD}$'';
``In Protocol 3, we start with $\rho_{BC}$ and act with $P_{B\to AB}$ \dots\ In the next step, we
act with the same map as in Protocol 2'';
``$\eta=\frac{L_A(L_C+L_D)}{(L_A+L_B)(L_B+L_C+L_D)}$'' (5.9);
``$-\log F^{(2)}\approx-\log F^{(3)}\propto f'c\eta^2$'' (5.10) ``with $f'<f$, confirming that the
sequential recovery is less effective than the single-step case''.

\emph{Use here.}  \S\ref{sec:vwz} decides 1(B) $=$ 1(A) and Protocol 3 exactly; the numerical
claim (5.10) is \emph{not} used (REF-P6-1 showed that Fig.~25 carries no information about the
corner kernel; the card \texttt{vwz-protocol-ordering} and track G3 own that question).

\smallskip\noindent\shot[0.85\textwidth]{G1a_VWZ_fig24.png}

\smallskip\noindent\shot[0.85\textwidth]{G1a_VWZ_510_511.png}
\end{citedbox}

\section{The strength of a corner: the exact first-order field}\label{sec:field}

\begin{lemma}[Admissibility]\label{lem:admissible}
In the notation of (C3), $k(M)\subseteq A_B$ if and only if $s\ge\lambda$, with equality
$k(M)=A_B$ iff $s=\lambda$.
\end{lemma}
\begin{proof}
\lstep{1}{1}{It suffices to treat the right case $k=\Rp_{p,\sigma}$, $M=(p,p+D)$,
 $D:=l_B+l_{\rm new}$, $A_B=(p,p+l_B)$.}
\lby{the left case is the image of the right one under the reflection $x\mapsto2p-x$, which maps
 $\Rp_{p,\sigma}$ to $\Lp_{p,\sigma}$ and preserves all lengths}
\lstep{1}{2}{$k$ is strictly increasing on $M$ with $k(p^+)=p$ and $k(p+D)=p+D/(1+s)$.}
\lby{\eqref{eq:step-map}, $k'(x)=(1+\sigma(x-p))^{-2}>0$, and $\sigma D=s$ by \eqref{eq:sigma}}
\lstep{1}{3}{$k(M)\subseteq A_B\iff D/(1+s)\le l_B\iff 1+s\ge D/l_B=1+\lambda\iff s\ge\lambda$.}
\lby{$\langle1\rangle2$, $D=l_B+l_{\rm new}$ and $\lambda=l_{\rm new}/l_B$}
\lqed{1}{4}
\end{proof}

\begin{definition}[Corner strength in a frame]\label{def:strength}
Let $p$ be a corner with parameter $\sigma$ and let $I'=(\alpha,\omega)$ be any interval with
$p\in I'$.  The \emph{strength of the corner in the frame $I'$} is
$\zeta^{(I')}:=\sigma\,\beta_{I'}(p)$, $\beta_{I'}(p)=(p-\alpha)(\omega-p)/(\omega-\alpha)$.
\end{definition}

\begin{lemma}[The parabolic parameter is a $(-1)$-density; exact]\label{lem:density}
Let $g\in\Mob$ be finite at $p$ (i.e.\ $g$ has no pole at $p$), and let
$h_\sigma(x)=p+(x-p)/(1+\sigma(x-p))$ be the parabolic flow fixing $p$.  Then
\begin{equation}\label{eq:density}
 g\circ h_\sigma\circ g^{-1}=\tilde h_{\sigma'}\ \text{ with }\ \sigma'=\frac\sigma{g'(p)},
\end{equation}
$\tilde h$ the parabolic flow fixing $g(p)$.  Consequently $\zeta^{(I')}$ of
Definition~\ref{def:strength} is invariant under $\Mob$: if $g\in\Mob$ is finite on
$\overline{I'}$ then the corner $(g(p),\sigma')$ in the frame $g(I')$ has the same strength.
\end{lemma}
\begin{proof}
\lstep{1}{1}{\ldefine $u:=1/(x-p)$ on $\{x\ne p\}$.  Then $h_\sigma$ reads $u\mapsto u+\sigma$.}
\lby{$1/(h_\sigma(x)-p)=(1+\sigma(x-p))/(x-p)=u+\sigma$}
\lstep{1}{2}{Write $g(x)=(ax+b)/(cx+d)$, $ad-bc>0$, $cp+d\ne0$.  Then
 $g(x)-g(p)=g'(p)(x-p)\bigl(1+\gamma(x-p)\bigr)^{-1}$ with $\gamma:=c/(cp+d)$.}
\lby{$g(x)-g(p)=\frac{(ad-bc)(x-p)}{(cx+d)(cp+d)}$, $g'(p)=\frac{ad-bc}{(cp+d)^2}$ and
 $cx+d=(cp+d)(1+\gamma(x-p))$}
\lstep{1}{3}{Hence $\tilde u:=1/(g(x)-g(p))=(u+\gamma)/g'(p)$, an affine function of $u$ with
 slope $1/g'(p)$.}
\lby{$\langle1\rangle2$ and $\langle1\rangle1$}
\lstep{1}{4}{Therefore the conjugate of $u\mapsto u+\sigma$ is
 $\tilde u\mapsto\tilde u+\sigma/g'(p)$, which is $\tilde h_{\sigma'}$ with
 $\sigma'=\sigma/g'(p)$: this is \eqref{eq:density}.}
\lby{$\langle1\rangle3$, $\langle1\rangle1$ applied at $g(p)$}
\lstep{1}{5}{Invariance of $\zeta$: $\beta_{g(I')}(g(p))=g'(p)\beta_{I'}(p)$ for $g\in\Mob$
 finite on $\overline{I'}$, so $\sigma'\beta_{g(I')}(g(p))=\sigma\beta_{I'}(p)$.}
\lby{$\beta_{I'}$ is the reciprocal of the derivative of the modular coordinate, and modular
 coordinates of $I'$ and $g(I')$ differ by the composition with $g$, so $\beta$ transforms as a
 vector field: $\beta_{g(I')}(g(x))=g'(x)\beta_{I'}(x)$ --- directly:
 $\frac{(g(x)-g(\alpha))(g(\omega)-g(x))}{g(\omega)-g(\alpha)}
 =g'(x)\frac{(x-\alpha)(\omega-x)}{\omega-\alpha}$ by $\langle1\rangle2$ applied three times;
 and $\langle1\rangle4$}
\lqed{1}{6}
\end{proof}

\begin{lemma}[Exact first-order field in a modular frame]\label{lem:field}
Let $I'=(\alpha,\omega)$ be an interval, $p\in I'$, $y=y_{I'}$ the modular coordinate
\eqref{eq:modframe} of $I'$ and $u:=y-y_{I'}(p)$.  Then, \emph{exactly} (not only to leading
order in $x-p$):
\begin{equation}\label{eq:field-identity}
 -\frac{(x-p)^2}{\beta_{I'}(x)}=-\beta_{I'}(p)\,\bigl(e^{u}-2+e^{-u}\bigr)
 =\beta_{I'}(p)\,w(u),\qquad w(u):=-4\sinh^2(u/2).
\end{equation}
Consequently the $\sigma$-derivative at $\sigma=0$ of the one-sided step \eqref{eq:step-map},
viewed as a vector field on $I'$ in the coordinate $y$, is
\begin{equation}\label{eq:field}
 \partial_\sigma\big|_{\sigma=0}\Rp_{p,\sigma}=\zeta\,w(u)\,\mathbf 1_{u>0}\,\partial_y,
 \qquad
 \partial_\sigma\big|_{\sigma=0}\Lp_{p,\sigma}=-\zeta\,w(u)\,\mathbf 1_{u<0}\,\partial_y,
 \qquad \zeta:=\sigma\beta_{I'}(p)\ \text{per unit }\sigma .
\end{equation}
For $I'=(-a,L)$, $p=0$ and $\sigma=s/L$ this gives $\zeta=as/(a+L)$, which is paper~1's
Definition of $\zeta$ (Box~[G2]).
\end{lemma}
\begin{proof}
\lstep{1}{1}{\ldefine $X:=(x-\alpha)/(\omega-x)$, $X_p:=(p-\alpha)/(\omega-p)$,
 $T':=\omega-\alpha$.  Then $e^u=X/X_p$ and
 $e^u-2+e^{-u}=(X-X_p)^2/(XX_p)$.}
\lby{\eqref{eq:modframe}: $y=\log X$, $y_{I'}(p)=\log X_p$; and
 $\frac XY+\frac YX-2=\frac{(X-Y)^2}{XY}$}
\lstep{1}{2}{$X-X_p=\dfrac{(x-p)T'}{(\omega-x)(\omega-p)}$.}
\lby{$(x-\alpha)(\omega-p)-(p-\alpha)(\omega-x)
 =\omega(x-p)-\alpha(x-p)=(x-p)T'$ after expanding both products}
\lstep{1}{3}{$XX_p=\dfrac{(x-\alpha)(p-\alpha)}{(\omega-x)(\omega-p)}$, hence
 $\dfrac{(X-X_p)^2}{XX_p}=\dfrac{(x-p)^2T'^2}{(\omega-x)(\omega-p)(x-\alpha)(p-\alpha)}$.}
\lby{$\langle1\rangle2$ and the definitions in $\langle1\rangle1$}
\lstep{1}{4}{$\beta_{I'}(p)\dfrac{(X-X_p)^2}{XX_p}
 =\dfrac{(x-p)^2T'}{(\omega-x)(x-\alpha)}=\dfrac{(x-p)^2}{\beta_{I'}(x)}$,
 which is \eqref{eq:field-identity}.}
\lby{$\beta_{I'}(p)=(p-\alpha)(\omega-p)/T'$, $\langle1\rangle3$, and
 $\beta_{I'}(x)=(x-\alpha)(\omega-x)/T'$}
\lstep{1}{5}{$e^u-2+e^{-u}=(e^{u/2}-e^{-u/2})^2=4\sinh^2(u/2)$, so
 $-\beta_{I'}(p)(e^u-2+e^{-u})=\beta_{I'}(p)w(u)$.}
\lby{the definition of $\sinh$}
\lstep{1}{6}{$\partial_\sigma|_{\sigma=0}\Rp_{p,\sigma}(x)=-(x-p)^2$ for $x>p$ and $0$ for
 $x\le p$; $\partial_\sigma|_{\sigma=0}\Lp_{p,\sigma}(x)=+(x-p)^2$ for $x<p$ and $0$ for
 $x\ge p$.}
\lby{differentiating \eqref{eq:step-map} in $\sigma$ at $\sigma=0$}
\lstep{1}{7}{A vector field $V(x)\partial_x$ reads $V(x)\beta_{I'}(x)^{-1}\partial_y$ in the
 coordinate $y$; with $\langle1\rangle6$, $\langle1\rangle4$ and $\langle1\rangle5$ this is
 \eqref{eq:field}, the indicator coming from $x\gtrless p\iff u\gtrless0$.}
\lby{$\partial_x=\frac{\dd y}{\dd x}\partial_y=\beta_{I'}(x)^{-1}\partial_y$ by
 \eqref{eq:modframe}}
\lstep{1}{8}{For $I'=(-a,L)$, $p=0$: $\beta_{I'}(0)=\frac{(0+a)(L-0)}{a+L}=\frac{aL}{a+L}$ and
 $\zeta=\frac sL\cdot\frac{aL}{a+L}=\frac{as}{a+L}$.}
\lby{Definition~\ref{def:strength} and $\sigma=s/L$}
\lqed{1}{9}
\end{proof}

\begin{remark}\label{rem:field-num}
\eqref{eq:field-identity} is an identity between rational functions and needs no smallness of
$x-p$; REF-P6-0 verified it in $50\,000$ random \emph{rational} configurations with $0$ failures
(\texttt{rigor/ref\_p6\_0\_field\_exact.out}, line~1), and the $(-1)$-density law
\eqref{eq:density} to $3\cdot10^{-6}$ in floating point (\emph{ibid.}, line~4).  The float test in
\texttt{ref\_p6\_0\_calculus.py}~(a) reports a spurious failure at relative error
$6.7\cdot10^{-3}$: it is pure cancellation as $x\to p$ and is superseded by the rational test
(REF-P6-0, ``Checks that PASSED'', item~1).
\end{remark}

\begin{remark}[Both signs of the compression give the same corner]\label{rem:signs}
The two fields in \eqref{eq:field} are supported on opposite sides of the corner and have
opposite signs; they are exchanged by $u\mapsto-u$ together with an overall sign, i.e.\ by the
reflection of $I'$.  Nevertheless both produce the \emph{same} Schwarzian mass $-2\sigma$ at $p$
(Lemma~\ref{lem:jump}), which is why the corner measure of (C4) does not record the side.
\end{remark}

\section{Schwarzian calculus for $C^1$ piecewise-M\"obius maps}\label{sec:schwarz}

\begin{definition}[The class $\mathrm{PM}(J)$]\label{def:PM}
Let $J\subseteq\mathbb R$ be an open interval.  $\mathrm{PM}(J)$ is the set of maps
$\Phi:J\to\mathbb R$ such that (i) $\Phi\in C^1(J)$ with $\Phi'>0$, and (ii) there are finitely
many \emph{breakpoints} $q_1<\dots<q_N$ in $J$ such that on the closure of each component of
$J\setminus\{q_1,\dots,q_N\}$, $\Phi$ coincides with (the restriction of) a M\"obius map without
pole there.  We write $v_\Phi:=\Phi''/\Phi'=(\log\Phi')'$, defined off the breakpoints, with
one-sided limits $v_\Phi(q^\pm)$ and jump $[v_\Phi]_q:=v_\Phi(q^+)-v_\Phi(q^-)$.
\end{definition}

\begin{lemma}[The Schwarzian of a $\mathrm{PM}$ map is a sum of point masses]\label{lem:masses}
Let $\Phi\in\mathrm{PM}(J)$ with breakpoints $q_1,\dots,q_N$.  Then $v_\Phi$ is of bounded
variation on compacts, $\Sch(\Phi)$ of \eqref{eq:schwarz} is a well-defined distribution, and
\begin{equation}\label{eq:masses}
 \Sch(\Phi)=\sum_{j=1}^N[v_\Phi]_{q_j}\,\delta_{q_j}.
\end{equation}
In particular $\Sch(\Phi)=0$ iff all jumps vanish, and there is no ``square'' contribution: the
term $-\frac12v^2$ carries no atom because $\Phi'$ is continuous.
\end{lemma}
\begin{proof}
\lstep{1}{1}{For a M\"obius $M(x)=(ax+b)/(cx+d)$ without pole on an interval,
 $v_M=-2c/(cx+d)$ and $v_M'-\frac12v_M^2=\frac{2c^2}{(cx+d)^2}-\frac{2c^2}{(cx+d)^2}=0$.}
\lby{$M'=(ad-bc)/(cx+d)^2$, $M''=-2c(ad-bc)/(cx+d)^3$, so $v_M=M''/M'=-2c/(cx+d)$}
\lstep{1}{2}{Hence on each component of $J\setminus\{q_j\}$ the function $v_\Phi$ is $C^\infty$,
 bounded on compacts of $\overline J$-components, and satisfies
 $v_\Phi'-\frac12v_\Phi^2=0$ pointwise.}
\lby{$\langle1\rangle1$ and Definition~\ref{def:PM}(ii)}
\lstep{1}{3}{$v_\Phi$ is a piecewise-$C^1$ function with finite one-sided limits at each $q_j$,
 hence of bounded variation on compacts, and its distributional derivative is
 $v_\Phi'=\{v_\Phi'\}+\sum_j[v_\Phi]_{q_j}\delta_{q_j}$, where $\{\cdot\}$ denotes the pointwise
 derivative off the breakpoints.}
\lby{the jump formula for the distributional derivative of a piecewise-$C^1$ function
 (integration by parts on each piece)}
\lstep{1}{4}{$v_\Phi^2\in L^\infty_{\rm loc}$, so it defines a distribution with no atoms, and
 $\Sch(\Phi)=\{v_\Phi'\}-\frac12v_\Phi^2+\sum_j[v_\Phi]_{q_j}\delta_{q_j}
 =\sum_j[v_\Phi]_{q_j}\delta_{q_j}$.}
\lby{$\langle1\rangle3$, $\langle1\rangle2$ and \eqref{eq:schwarz}}
\lstep{1}{5}{The parenthetical claim: $\Phi'''$ as a distribution equals
 $\{\Phi'''\}+\sum_j[\Phi'']_{q_j}\delta_{q_j}$ and $[\Phi'']_{q_j}=\Phi'(q_j)[v_\Phi]_{q_j}$
 because $\Phi'$ is continuous at $q_j$; dividing by the continuous positive function $\Phi'$ is
 legitimate on an atom and reproduces \eqref{eq:masses}.}
\lby{$\Phi\in C^1$ (Definition~\ref{def:PM}(i)) and $\Phi''=v_\Phi\Phi'$}
\lqed{1}{6}
\end{proof}

\begin{lemma}[The mass of a parabolic corner: $-2\sigma$ from either side]\label{lem:jump}
For $\sigma>0$, $\Rp_{p,\sigma}$ and $\Lp_{p,\sigma}$ of \eqref{eq:step-map} belong to
$\mathrm{PM}$ near $p$, with the single breakpoint $p$, and
\begin{equation}\label{eq:jump}
 [v]_p=-2\sigma\qquad\text{for both.}
\end{equation}
\end{lemma}
\begin{proof}
\lstep{1}{1}{$\Rp_{p,\sigma}$ is $C^1$ at $p$: the left derivative is $1$ and the right
 derivative is $\lim_{x\downarrow p}(1+\sigma(x-p))^{-2}=1$.}
\lby{\eqref{eq:step-map} and $\frac{\dd}{\dd x}\frac{x-p}{1+\sigma(x-p)}=(1+\sigma(x-p))^{-2}$}
\lstep{1}{2}{For $x>p$: $v(x)=-2\sigma/(1+\sigma(x-p))$, so $v(p^+)=-2\sigma$; for $x<p$:
 $v(x)=0$, so $v(p^-)=0$.  Hence $[v]_p=-2\sigma$.}
\lby{$\Phi'=(1+\sigma(x-p))^{-2}$, $\Phi''=-2\sigma(1+\sigma(x-p))^{-3}$ for $x>p$; $\Phi=\id$
 for $x<p$}
\lstep{1}{3}{For $\Lp_{p,\sigma}$ and $x<p$: $\Phi'=(1-\sigma(x-p))^{-2}$,
 $\Phi''=+2\sigma(1-\sigma(x-p))^{-3}$, so $v(x)=2\sigma/(1-\sigma(x-p))$ and $v(p^-)=2\sigma$;
 $v(p^+)=0$.  Hence $[v]_p=0-2\sigma=-2\sigma$.}
\lby{differentiating \eqref{eq:step-map} (left case) twice}
\lqed{1}{4}
\end{proof}

\begin{lemma}[Composition]\label{lem:compose}
Let $g\in\mathrm{PM}(J)$ and $f\in\mathrm{PM}(g(J))$.  Then $f\circ g\in\mathrm{PM}(J)$, its
breakpoint set is contained in $\{$breakpoints of $g\}\cup g^{-1}\{$breakpoints of $f\}$, and
\begin{equation}\label{eq:comp-rule}
 \Sch(f\circ g)=\Sch(g)+\sum_{r}m^f_r\,g'(g^{-1}(r))\,\delta_{g^{-1}(r)},
 \qquad \Sch(f)=\sum_rm^f_r\delta_r,
\end{equation}
i.e.\ $\Sch(f\circ g)=\Sch(g)+(g')^2\,\Sch(f)\circ g$ with the pullback of a point mass
$\delta_r\mapsto g'(g^{-1}(r))^{-1}\delta_{g^{-1}(r)}$.  Also $g^{-1}\in\mathrm{PM}(g(J))$.
\end{lemma}
\begin{proof}
\lstep{1}{1}{$f\circ g$ is $C^1$ with $(f\circ g)'=(f'\circ g)g'>0$, and on each component of
 $J$ minus the stated finite set both $f$ and $g$ are M\"obius, hence so is $f\circ g$: thus
 $f\circ g\in\mathrm{PM}(J)$ with the stated breakpoints.}
\lby{the chain rule, $f,g\in C^1$ with positive derivatives, and the fact that $\Mob$ is a group}
\lstep{1}{2}{Off the breakpoints, $v_{f\circ g}=(v_f\circ g)\,g'+v_g$.}
\lby{$\log(f\circ g)'=\log(f'\circ g)+\log g'$, differentiated}
\lstep{1}{3}{At a point $x_0$ of the finite set, $[v_{f\circ g}]_{x_0}
 =g'(x_0)\,[v_f]_{g(x_0)}+[v_g]_{x_0}$.}
\lby{$\langle1\rangle2$, continuity of $g$ and of $g'$ at $x_0$ (Definition~\ref{def:PM}(i)), and
 the fact that one-sided limits of a product of a continuous and a one-sided-continuous factor
 multiply; $g$ increasing, so the sides are not exchanged}
\lstep{1}{4}{Summing $\langle1\rangle3$ over the finite set and using \eqref{eq:masses} three
 times gives \eqref{eq:comp-rule}.}
\lby{Lemma~\ref{lem:masses} for $f\circ g$, for $g$ and for $f$; $[v_f]_{g(x_0)}=m^f_{g(x_0)}$}
\lstep{1}{5}{$g^{-1}\in\mathrm{PM}(g(J))$.}
\lby{$g$ is a $C^1$ increasing bijection onto $g(J)$ with $g'>0$, so $g^{-1}\in C^1$ with
 $(g^{-1})'=1/(g'\circ g^{-1})>0$; on each piece $g$ is M\"obius, hence so is $g^{-1}$}
\lqed{1}{6}
\end{proof}

\begin{lemma}[Two maps with the same Schwarzian differ by a global M\"obius
map]\label{lem:same-schwarz}
Let $\Phi_1,\Phi_2\in\mathrm{PM}(J)$.  Then $\Sch(\Phi_1)=\Sch(\Phi_2)$ if and only if there is
$g\in\Mob$, without pole on the \emph{open} interval $\Phi_1(J)$, with $\Phi_2=g\circ\Phi_1$.
(The closure may not be excluded: $\Phi_1=\id$, $\Phi_2(x)=x/(1-x)$ on $J=(0,1)$ have the same
Schwarzian and force $g=\Phi_2$, whose pole is at $1\in\overline{\Phi_1(J)}$; REF-P6-3 \S5.)
\end{lemma}
\begin{proof}
\lstep{1}{1}{($\Leftarrow$) If $\Phi_2=g\circ\Phi_1$ with $g\in\Mob$, then $\Sch(g)=0$ and
 \eqref{eq:comp-rule} (with $f=g$, $g=\Phi_1$) gives $\Sch(\Phi_2)=\Sch(\Phi_1)$.}
\lby{Lemma~\ref{lem:masses} step $\langle1\rangle1$ ($g$ has no breakpoint) and
 Lemma~\ref{lem:compose}}
\lstep{1}{2}{($\Rightarrow$) \ldefine $\Psi:=\Phi_2\circ\Phi_1^{-1}\in\mathrm{PM}(\Phi_1(J))$.
 Then $\Phi_2=\Psi\circ\Phi_1$ and, by \eqref{eq:comp-rule},
 $0=\Sch(\Phi_2)-\Sch(\Phi_1)=\sum_rm^\Psi_r\Phi_1'(\Phi_1^{-1}(r))\delta_{\Phi_1^{-1}(r)}$.}
\lby{Lemma~\ref{lem:compose} ($\Phi_1^{-1}\in\mathrm{PM}$, composition of $\mathrm{PM}$ maps) and
 the hypothesis}
\lstep{1}{3}{Hence $m^\Psi_r=0$ for every $r$, i.e.\ all jumps of $v_\Psi$ vanish and $v_\Psi$ is
 continuous on $\Phi_1(J)$.}
\lby{$\Phi_1'>0$ and the fact that point masses at distinct points are linearly independent}
\lstep{1}{4}{Let $q$ be a breakpoint of $\Psi$ with M\"obius pieces $M_-$ (left) and $M_+$
 (right).  Then $M_-(q)=M_+(q)$, $M_-'(q)=M_+'(q)$ and $M_-''(q)=M_+''(q)$.}
\lby{$\Psi\in C^1$ gives the first two; $\langle1\rangle3$ and $\Psi''=v_\Psi\Psi'$ give the
 third}
\lstep{1}{5}{A M\"obius map without pole at $q$ is determined by
 $(M(q),M'(q),M''(q))$; hence $M_-=M_+$ and $q$ is not a breakpoint after all.}
\lby{Lemma~\ref{lem:density} $\langle1\rangle2$: $M(x)=M(q)+M'(q)(x-q)/(1+\gamma(x-q))$ with
 $\gamma=-M''(q)/(2M'(q))$, so the three data determine $M$}
\lstep{1}{6}{Therefore $\Psi$ has no breakpoints and is a single M\"obius map $g$ on
 $\Phi_1(J)$, without pole there, and $\Phi_2=g\circ\Phi_1$.}
\lby{$\langle1\rangle5$ applied to each breakpoint in turn, and Definition~\ref{def:PM}(ii)}
\lqed{1}{7}
\end{proof}

\begin{verdictok}
Lemmas~\ref{lem:masses}--\ref{lem:same-schwarz} are the three Schwarzian facts asserted in
\texttt{rigor/phase6\_brief.md} (\S``SCHWARZIAN''), and they hold as stated: masses
$[\Phi''/\Phi'](p_k)$ with no square term, $-2\sigma$ for a parabolic corner \emph{from either
side}, the composition rule, and M\"obius rigidity.  Numerics: \texttt{g1a\_calculus.py} \S1
(jumps of $\Lp\circ\Rp$ and $\Rp\circ\Lp$ equal $-2(\sigma_1+\sigma_2)$ to $10^{-4}$) and
REF-P6-0 items 2--3.
\end{verdictok}

\section{The corner calculus}\label{sec:corner}

\begin{definition}[Hypothesis (H)]\label{def:H}
A one-sided protocol $\PP$ (C4) satisfies \emph{hypothesis (H)} if
\begin{equation}\label{eq:H}
 p^{(m)}\notin\operatorname{int}M_{m'}\qquad\text{for all }1\le m<m'\le N,
\end{equation}
i.e.\ no \emph{earlier} corner lies in the interior of a \emph{later} step's moving part.
\end{definition}

\begin{theorem}[General transformation rule; no hypothesis]\label{thm:general}
Let $\PP$ be any one-sided protocol and $\Phi=k_1\circ\dots\circ k_N$ its composite point map.
Put $G_m:=k_{m+1}\circ\dots\circ k_N$ (so $G_N=\id$) and, for those $m$ with
$p^{(m)}\in G_m(J_N)$, $x_m:=G_m^{-1}(p^{(m)})$.  Then $\Phi\in\mathrm{PM}(J_N)$ and
\begin{equation}\label{eq:general-rule}
 \Sch(\Phi)=-2\sum_{m:\,p^{(m)}\in G_m(J_N)}\sigma_m\,G_m'(x_m)\,\delta_{x_m},
\end{equation}
masses at coinciding $x_m$ being added.  Corners $m$ with $p^{(m)}\notin G_m(J_N)$ contribute
nothing at all.
\end{theorem}
\begin{proof}
\lstep{1}{1}{Each $k_m$ belongs to $\mathrm{PM}(J_m)$, with the single breakpoint $p^{(m)}$ and
 $\Sch(k_m)=-2\sigma_m\delta_{p^{(m)}}$.}
\lproof
\lstep{2}{1}{$k_m$ is the identity on one side of $p^{(m)}$ and equals the M\"obius map
 $x\mapsto p+(x-p)/(1\pm\sigma_m(x-p))$ on the other.}
\lby{(C3), \eqref{eq:step-map}}
\lstep{2}{2}{That M\"obius map has no pole on the closure of the moving part: the pole is at
 $p\mp1/\sigma_m$, which lies on the \emph{identity} side of $p$.}
\lby{for $\Rp$: pole at $p-1/\sigma_m<p$ while the piece is $x\ge p$; for $\Lp$: pole at
 $p+1/\sigma_m>p$ while the piece is $x\le p$}
\lstep{2}{3}{$k_m\in C^1$ with $k_m'>0$, and the jump is $-2\sigma_m$.}
\lby{Lemma~\ref{lem:jump}}
\lqed{2}{4} \lby{$\langle2\rangle1$--$\langle2\rangle3$ and Definition~\ref{def:PM},
 Lemma~\ref{lem:masses}}
\lstep{1}{2}{$\Phi\in\mathrm{PM}(J_N)$ and $\Sch(\Phi)=\Sch(G_1)+\sum_r m^{k_1}_r
 G_1'(G_1^{-1}(r))\delta_{G_1^{-1}(r)}$, the sum running over $r\in G_1(J_N)$.}
\lby{$\Phi=k_1\circ G_1$, $\langle1\rangle1$, Lemma~\ref{lem:compose}; the composition rule
 \eqref{eq:comp-rule} only sees breakpoints of $k_1$ that lie in the range $G_1(J_N)$, since
 $f=k_1$ is restricted to that range}
\lstep{1}{3}{Iterating $\langle1\rangle2$ over $m=1,2,\dots,N$ gives \eqref{eq:general-rule}.}
\lby{induction on $N$: $G_m=k_{m+1}\circ G_{m+1}$, $\Sch(k_m)=-2\sigma_m\delta_{p^{(m)}}$ by
 $\langle1\rangle1$, and the chain rule $(k_{m+1}\circ\dots\circ k_N)'=G_m'$}
\lqed{1}{4}
\end{proof}

\begin{theorem}[Corner calculus under (H)]\label{thm:corner}
Let $\PP$ be a one-sided protocol satisfying (H).  Then every corner $p^{(m)}$ is one of the
points $p_1,\dots,p_{n+1}$, $\Phi\in\mathrm{PM}(J_N)$, and
\begin{equation}\label{eq:corner-calculus}
 \boxed{\ \Sch(\Phi)=-2\sum_{k:\,p_k\in\operatorname{int}J_N}\kappa_k\,\delta_{p_k},\qquad
 \kappa_k=\sum_{m:\,p^{(m)}=p_k}\sigma_m\ }
\end{equation}
--- the corner measure of (C4).  In particular $\Sch(\Phi)$ is independent of the order in which
the steps are performed, as long as every order considered is a legal protocol satisfying (H)
with the same multiset of (corner, parameter) pairs.
\end{theorem}
\begin{proof}
\lstep{1}{1}{Every corner is a chain point: $p^{(m)}$ is an endpoint of the block
 $A^{(m)}_B$, and endpoints of blocks are among $p_1,\dots,p_{n+1}$.}
\lby{(C1), (C3)}
\lstep{1}{2}{\lassume $m<m'$.  \lprove $k_{m'}(p^{(m)})=p^{(m)}$ and $k_{m'}'(p^{(m)})=1$.}
\lproof
\lstep{2}{1}{$p^{(m)}\in\overline{J_{m-1}}\subseteq\overline{J_{m'-1}}$.}
\lby{$p^{(m)}$ is an endpoint of $A^{(m)}_B\subseteq J_{m-1}$, and $J_0\subseteq J_1\subseteq
 \dots$ by (C4)}
\lstep{2}{2}{Let $f_{m'}$ be the endpoint of $M_{m'}$ other than $p^{(m')}$ (the outer end of
 $A^{(m')}_{\rm new}$).  Then $f_{m'}\notin\overline{J_{m'-1}}$, hence $p^{(m)}\ne f_{m'}$.}
\lby{$A^{(m')}_{\rm new}$ is disjoint from $J_{m'-1}$ and shares exactly one endpoint $q$ with
 it; $f_{m'}$ is at distance $l^{(m')}_{\rm new}>0$ from $q$ on the far side; and
 $\langle2\rangle1$}
\lstep{2}{3}{Hence either $p^{(m)}\notin\overline{M_{m'}}$, or $p^{(m)}=p^{(m')}$.}
\lby{$\overline{M_{m'}}=\operatorname{int}M_{m'}\cup\{p^{(m')},f_{m'}\}$, (H) excludes the
 interior, and $\langle2\rangle2$ excludes $f_{m'}$}
\lstep{2}{4}{\lcase $p^{(m)}\notin\overline{M_{m'}}$: $k_{m'}$ is the identity on a neighbourhood
 of $p^{(m)}$, so $k_{m'}(p^{(m)})=p^{(m)}$ and $k_{m'}'(p^{(m)})=1$.}
\lby{(C3): $k_{m'}=\id$ off $M_{m'}$, and $\overline{M_{m'}}$ is closed}
\lstep{2}{5}{\lcase $p^{(m)}=p^{(m')}$: $k_{m'}(p^{(m')})=p^{(m')}$ and $k_{m'}'(p^{(m')})=1$.}
\lby{\eqref{eq:step-map} fixes $p$, and Lemma~\ref{lem:jump} $\langle1\rangle1$ (the map is $C^1$
 at $p$ with derivative $1$ on both sides)}
\lqed{2}{6} \lby{$\langle2\rangle3$--$\langle2\rangle5$}
\lstep{1}{3}{$G_m(p^{(m)})=p^{(m)}$ and $G_m'(p^{(m)})=1$; in particular
 $p^{(m)}\in G_m(J_N)$ and $x_m=p^{(m)}$.}
\lby{$G_m=k_{m+1}\circ\dots\circ k_N$, $\langle1\rangle2$ applied to each factor in turn, and the
 chain rule (a product of factors all equal to $1$); $p^{(m)}\in J_N$ since $J_N\supseteq J_{m-1}$
 and $p^{(m)}\in\overline{J_{m-1}}$ --- corners at $\partial J_N$ (which occur only for a
 degenerate step $A_B^{(m)}=J_{m-1}$, Remark~\ref{rem:degenerate}) carry no mass in
 $\Sch(\Phi)|_{J_N}$, which is why the sum in \eqref{eq:corner-calculus} runs over the interior
 junctions}
\lstep{1}{4}{Substituting $\langle1\rangle3$ into \eqref{eq:general-rule} gives
 $\Sch(\Phi)=-2\sum_m\sigma_m\delta_{p^{(m)}}$, which is \eqref{eq:corner-calculus} after
 grouping the $m$'s by the value of $p^{(m)}$ (a finite set of chain points by
 $\langle1\rangle1$).}
\lby{Theorem~\ref{thm:general} and $\langle1\rangle3$}
\lstep{1}{5}{Order independence: \eqref{eq:corner-calculus} depends only on the multiset
 $\{(p^{(m)},\sigma_m)\}$.}
\lby{$\langle1\rangle4$}
\lqed{1}{6}
\end{proof}

\begin{proposition}[Sufficient conditions for (H)]\label{prop:H-suff}
Each of the following implies (H).
\begin{enumerate}[label=(\alph*),leftmargin=2.2em]
\item Every conditioning region $A^{(m)}_B$ is a \emph{single chain interval} \emph{and} the
 starting block $J_0$ consists of at least two chain intervals.
\item All corners coincide: $p^{(1)}=\dots=p^{(N)}$.
\item $N=1$.
\end{enumerate}
Condition (a) is \emph{not} implied by ``every $A^{(m)}_B$ is a single chain interval'' alone
(Example~\ref{ex:swallow}), and (H) does not imply that the $A_B^{(m)}$ are single intervals
(VWZ Protocol 1(B), \S\ref{sec:vwz}, satisfies (b) but conditions on a union).
\end{proposition}
\begin{proof}
\lstep{1}{1}{(c) is vacuous and (b) is immediate: $p^{(m)}=p^{(m')}$ is an \emph{endpoint} of
 $M_{m'}$, not an interior point.}
\lby{Definition~\ref{def:H} and \eqref{eq:corner-def}}
\lstep{1}{2}{\lassume (a) and $m<m'$.  \lprove $p^{(m)}\notin\operatorname{int}M_{m'}$.}
\lproof
\lstep{2}{1}{$J_{m-1}$ consists of at least two chain intervals and $A^{(m)}_B$ is a single one
 of them, so $A_B^{(m)}\subsetneq J_{m-1}$ and its inner endpoint $p^{(m)}$ is an
 \emph{interior} point of $J_{m-1}$, hence of $J_{m'-1}$.}
\lby{(C3): $p^{(m)}$ is the endpoint of $A^{(m)}_B$ shared with $J_{m-1}\setminus A^{(m)}_B\ne
 \emptyset$; $J_0\subseteq J_{m-1}\subseteq J_{m'-1}$}
\lstep{2}{2}{$\operatorname{int}M_{m'}$ contains exactly one chain point, namely the junction
 $q_{m'}$ between $A^{(m')}_B$ and $A^{(m')}_{\rm new}$, and $q_{m'}$ is an endpoint of
 $J_{m'-1}$.}
\lby{$M_{m'}$ is the union of the two adjacent chain intervals $A^{(m')}_B$ and
 $A^{(m')}_{\rm new}$ (both single), so its interior meets the chain points only in their common
 junction; $q_{m'}\in\partial J_{m'-1}$ because $A^{(m')}_{\rm new}$ is attached to $J_{m'-1}$
 there}
\lstep{2}{3}{$p^{(m)}$ is a chain point (Theorem~\ref{thm:corner} $\langle1\rangle1$) and
 $p^{(m)}\ne q_{m'}$, because $p^{(m)}\in\operatorname{int}J_{m'-1}$ and
 $q_{m'}\in\partial J_{m'-1}$.}
\lby{$\langle2\rangle1$, $\langle2\rangle2$, and $\operatorname{int}J\cap\partial J=\emptyset$}
\lqed{2}{4} \lby{$\langle2\rangle2$ and $\langle2\rangle3$: the only chain point in
 $\operatorname{int}M_{m'}$ is $q_{m'}\ne p^{(m)}$}
\lqed{1}{3}
\end{proof}

\subsection{Two failures of (H), and what they cost}\label{sec:failH}

\begin{example}[REF-P6-0: union conditioning displaces the mass off the
junction]\label{ex:displaced}
Chain $n=4$, $(l_1,l_2,l_3,l_4)=(1.3,2.1,0.9,1.7)$, so $(p_1,\dots,p_5)=(0,1.3,3.4,4.3,6.0)$;
$J_0=A_2A_3$; ordinary Petz steps.  Step~1 adjoins $A_4$ conditioning on $A_3$: corner $p_3=3.4$,
$\sigma_1=2l_4/(l_3(l_3+l_4))=1.452991$, $k_1=\Rp_{p_3,\sigma_1}$.  Step~2 adjoins $A_1$
conditioning on the \emph{union} $A_2A_3$: corner $p_4=4.3$,
$\sigma_2=2l_1/((l_2+l_3)(l_1+l_2+l_3))=0.201550$, $k_2=\Lp_{p_4,\sigma_2}$ (admissible: the
image of $p_1$ is $1.9964\ge p_2$).  Here $p^{(1)}=p_3\in\operatorname{int}M_2=(p_1,p_4)$, so (H)
fails, and Theorem~\ref{thm:general} gives a mass
$-2\sigma_1k_2'(x_0)=-1.947339$ at $x_0=k_2^{-1}(p_3)=3.200568$ ($k_2'(x_0)=0.670114$), no mass at
$p_3$, and $-2\sigma_2$ at $p_4$; the total mass is $71\%$ of the naive
$-2(\sigma_1+\sigma_2)$.  Measured: $-1.947332$ at $x_0$, $+5.4\cdot10^{-5}$ at $p_3$
(\texttt{ref\_p6\_0\_w1\_counterexample.out}; reproduced in \texttt{g1a\_calculus.py} \S2).
Control: with $A_B=A_2$ (a single interval) both corners sit at $p_3$ and the masses add
($-3.634266$ against $-2(\sigma_1+\sigma_2)=-3.634274$).
\end{example}

\begin{example}[New: single-interval conditioning is not enough --- a corner can be
\emph{swallowed}]\label{ex:swallow}
Chain $n=3$, $l=(1,1,1)$, $(p_1,\dots,p_4)=(0,1,2,3)$, ordinary Petz, starting block
$J_0=A_2$ --- a \emph{single} chain interval.  Step~1 adjoins $A_3$ conditioning on $A_2$; since
$J_0=A^{(1)}_B$, the corner is the far end $p^{(1)}=p_2=1$ and
$\sigma_1=2l_3/(l_2(l_2+l_3))=1$, $k_1=\Rp_{p_2,1}$.  Step~2 adjoins $A_1$ conditioning on
$A^{(2)}_B=A_2$, a single chain interval; now $J_1=A_2A_3$, so the corner is the inner end
$p^{(2)}=p_3=2$, $\sigma_2=1$, $k_2=\Lp_{p_3,1}$.  Both conditioning regions are single chain
intervals, yet $p^{(1)}=p_2=1\in\operatorname{int}M_2=(p_1,p_3)=(0,2)$: (H) fails.  Moreover
$k_2\bigl((p_1,p_3)\bigr)=(4/3,2)$ does not contain $p_2$, so in Theorem~\ref{thm:general} the
index $m=1$ does not contribute \emph{at all}:
\[
 \Sch(\Phi)=-2\sigma_2\,\delta_{p_3},\qquad \Phi=g\circ k_2\ \text{ with } g\in\Mob,
\]
i.e.\ the entire first step is undone by a global M\"obius map and leaves the recovered state
unchanged (Theorem~\ref{thm:state}(3)).  Verified in \texttt{g1a\_calculus.py} \S3
(\texttt{g1a\_calculus.out}).
\end{example}

\begin{verdictgap}
\textbf{Correction to the wording of W1 / card \texttt{corner-calculus} (MINOR).}  Both the brief
and the card say ``every step's conditioning region $A_B$ is a single interval of the chain, or,
more generally, no earlier corner lies in the interior of a later step's moving part
\dots\ the first implies the second''.  Example~\ref{ex:swallow} shows that the implication
\emph{fails} when the starting block is a single chain interval: there both $A_B$'s are single
chain intervals and (H) is violated.  The correct sufficient condition is
Proposition~\ref{prop:H-suff}(a): single-interval conditioning \emph{and} a starting block with at
least two chain intervals.  Nothing downstream is affected: VWZ's protocols and the tables of
\S\ref{sec:tables} all start from an adjacent pair, and the brief's own protocol class is
``start from any adjacent pair''.  The weaker condition (H) remains the one the proof uses, and
Theorem~\ref{thm:general} covers every case.
\end{verdictgap}

\begin{verdictok}
Theorem~\ref{thm:corner} is W1 as repaired by REF-P6-0/1/2, now \emph{proved}, together with the
unconditional Theorem~\ref{thm:general}.  The point-order argument is exactly the one sketched in
the brief; the two ingredients that make it work are isolated as
Theorem~\ref{thm:corner} $\langle1\rangle2$ (later steps fix earlier corners with derivative $1$)
and Lemma~\ref{lem:compose} (the pullback factor $G_m'(x_m)$).  Numerics:
\texttt{g1a\_calculus.py} \S\S1--5, all $33$ checks PASS (\texttt{g1a\_calculus.out}).
\end{verdictok}

\section{The recovered state depends only on the Schwarzian}\label{sec:state}

\begin{lemma}[Functoriality of the half-density push-forward]\label{lem:comp-W}
Let $k_2:J_2\to J_1$ and $k_1:J_1\to J_0$ be orientation-preserving $C^1$ diffeomorphisms onto
their ranges.  Then $V_{k_1}V_{k_2}=V_{k_1\circ k_2}$ and each $V_k$ is an isometry
$H_J\to H_{k(J)}$; if $k$ is a bijection of $\mathbb R$ then $V_k$ is unitary on $H$.
\end{lemma}
\begin{proof}
\lstep{1}{1}{$\lVert V_kf\rVert^2=\int_{k(J)}|(k^{-1})'(y)||f(k^{-1}(y))|^2\dd y
 =\int_J|f(x)|^2\dd x$.}
\lby{\eqref{eq:pushforward} and the substitution $y=k(x)$, $\dd y=k'(x)\dd x$, $(k^{-1})'(k(x))
 =1/k'(x)$}
\lstep{1}{2}{$(V_{k_1}V_{k_2}f)(y)=\bigl((k_1^{-1})'(y)\bigr)^{1/2}
 \bigl((k_2^{-1})'(k_1^{-1}(y))\bigr)^{1/2}f\bigl(k_2^{-1}k_1^{-1}(y)\bigr)
 =\bigl(((k_1k_2)^{-1})'(y)\bigr)^{1/2}f\bigl((k_1k_2)^{-1}(y)\bigr)$.}
\lby{\eqref{eq:pushforward} twice and the chain rule
 $((k_1\circ k_2)^{-1})'=(k_2^{-1})'\circ k_1^{-1}\cdot(k_1^{-1})'$; all factors are positive, so
 the positive square roots multiply}
\lqed{1}{3}
\end{proof}

\begin{lemma}[M\"obius invariance of the vacuum, one-particle level]\label{lem:moebius-inv}
Let $g\in\Mob$, i.e.\ $g(x)=(ax+b)/(cx+d)$ with $a,b,c,d$ real and $ad-bc>0$.  Then $V_g$ is
unitary on $H=L^2(\mathbb R;\mathbb C^r)$ and $[V_g,P]=0$.  Consequently
$\omega\circ\beta_{V_g}=\omega$ on $\Fr(\mathbb R)$, and $\omega\circ\beta_{V_g}=\omega$ on
$\Fr(J)$ for every interval $J$ on which $g$ has no pole.
\end{lemma}
\begin{proof}
\lstep{1}{1}{$g$ is a bijection of $\mathbb R\cup\{\infty\}$ and $g'>0$ off its pole, so $V_g$ is
 unitary on $H$ (a single point is a Lebesgue null set).}
\lby{$ad-bc>0$ gives $g'=(ad-bc)/(cx+d)^2>0$; Lemma~\ref{lem:comp-W} $\langle1\rangle1$ and
 surjectivity}
\lstep{1}{2}{$P$ is the projection onto the Hardy space $H^2$ of boundary values of functions
 holomorphic and $L^2$-bounded in the upper half-plane $\mathbb C_+$.}
\lby{(C2); Longo--Xu \cite{LX} \S3.1 (Hardy projection), \texttt{cited\_R1.tex} Box~C1}
\lstep{1}{3}{$g$ extends to a biholomorphism of $\mathbb C_+$ onto itself, and
 $(g^{-1})'(w)=(ad-bc)/(a-cw)^2$ is holomorphic and zero-free on $\mathbb C_+$, so it has a
 single-valued holomorphic square root there (take $\sqrt{ad-bc}/(a-cw)$, whose boundary values
 on $\mathbb R$ are positive up to a constant phase).}
\lby{$\mathrm{PSL}(2,\mathbb R)$ acts on $\mathbb C_+$; $g^{-1}(w)=(dw-b)/(-cw+a)$ and $a-cw\ne0$
 for $w\in\mathbb C_+$ because $a,c$ are real}
\lstep{1}{4}{Hence $f\in H^2\Rightarrow V_gf\in H^2$: $V_gf$ is the boundary value of
 $w\mapsto\bigl((g^{-1})'(w)\bigr)^{1/2}f(g^{-1}(w))$, holomorphic on $\mathbb C_+$, and
 $\lVert V_gf\rVert_2=\lVert f\rVert_2<\infty$.}
\lby{$\langle1\rangle3$, $\langle1\rangle2$, $\langle1\rangle1$}
\lstep{1}{5}{Applying $\langle1\rangle4$ to $g$ and to $g^{-1}$ gives $V_gH^2=H^2$, hence
 $V_gP=PV_g$.}
\lby{$V_{g^{-1}}=V_g^{-1}=V_g^*$ by Lemma~\ref{lem:comp-W}, so $V_gH^2\subseteq H^2$ and
 $V_g^*H^2\subseteq H^2$; a unitary with $V\cP\subseteq\cP$ and $V^*\cP\subseteq\cP$ commutes with
 the projection onto $\cP$}
\lstep{1}{6}{$\omega\circ\beta_{V_g}=\omega$: the two-point function of $\omega\circ\beta_{V_g}$
 is $\langle V_gf,PV_gh\rangle=\langle f,V_g^*PV_gh\rangle=\langle f,Ph\rangle$, and a quasi-free
 state is determined by its two-point function.}
\lby{$\langle1\rangle5$, $V_g$ unitary, and the definition of a gauge-invariant quasi-free state
 (C2)}
\lqed{1}{7}
\end{proof}

\begin{theorem}[Normality, functoriality, and Schwarzian dependence]\label{thm:state}
Let $\PP$ be a one-sided protocol with composite point map $\Phi$ and put
$\beta_\PP:=\beta_1\circ\beta_2\circ\dots\circ\beta_N$, $\beta_m:=\overline\beta_{s_m}$ the
zero-collar step maps of Box~[G1].  Then:
\begin{enumerate}[label=(\arabic*),leftmargin=2.2em]
\item $\beta_\PP:\Fr(J_N)\to\Fr(J_0)$ is a normal, parity-preserving $*$-isomorphism onto its
 range, and $\beta_\PP=\beta_{V_\Phi}$ on the CAR generators;
\item if $\Phi_2=g\circ\Phi_1$ with $g\in\Mob$ (no condition on the pole of $g$ is needed: by
 Lemma~\ref{lem:moebius-inv} $V_g$ is unitary on all of $H$), then
 $\omega\circ\beta_{V_{\Phi_2}}=\omega\circ\beta_{V_{\Phi_1}}$;
\item consequently, by Lemma~\ref{lem:same-schwarz}, the recovered state
 $\omega\circ\beta_\PP$ on $\Fr(J_N)$ depends on $\PP$ \emph{only through} $\Sch(\Phi)$; under
 hypothesis (H), only through the corner measure $\kappa_\PP$.
\end{enumerate}
\end{theorem}
\begin{proof}
\lstep{1}{1}{Each $\beta_m$ is a normal $*$-isomorphism of $\Fr(J_m)$ onto its range inside
 $\Fr(J_{m-1})$, and on CAR generators $\beta_m=\beta_{V_{k_m}}$.}
\lby{Box~[G1] (Note~4 Thm.~7.1) applied to $A=J_{m-1}\setminus A^{(m)}_B$,
 $B=A^{(m)}_B$, $C=A^{(m)}_{\rm new}$, whose hypothesis $s_m\ge\lambda_m$ holds by (C3) and
 Lemma~\ref{lem:admissible}; the $C^*$-level map $\beta_{s}$ there is by construction the
 half-density push-forward along the geometric compression, i.e.\ $\beta_{V_{k_m}}$}
\lstep{1}{2}{A composition of normal $*$-isomorphisms onto their ranges is a normal
 $*$-isomorphism onto its range.}
\lby{normality (\(\sigma\)-weak continuity) and the isomorphism property are both preserved under
 composition; the range of $\beta_m$ is a von Neumann subalgebra of $\Fr(J_{m-1})$, on which
 $\beta_{m-1}$ is defined and normal}
\lstep{1}{3}{$\beta_\PP=\beta_{V_{k_1}}\circ\dots\circ\beta_{V_{k_N}}
 =\beta_{V_{k_1}\cdots V_{k_N}}=\beta_{V_\Phi}$ on CAR generators.}
\lby{$\langle1\rangle1$, Bogoliubov functoriality $\beta_{V_1}\beta_{V_2}=\beta_{V_1V_2}$ (C2),
 Lemma~\ref{lem:comp-W} and \eqref{eq:Phi}}
\lstep{1}{4}{(2): $V_{\Phi_2}=V_gV_{\Phi_1}$, hence
 $\omega\circ\beta_{V_{\Phi_2}}=(\omega\circ\beta_{V_g})\circ\beta_{V_{\Phi_1}}
 =\omega\circ\beta_{V_{\Phi_1}}$.}
\lby{Lemma~\ref{lem:comp-W}, Bogoliubov functoriality, and Lemma~\ref{lem:moebius-inv}}
\lstep{1}{5}{(3): if $\Sch(\Phi_1)=\Sch(\Phi_2)$ then $\Phi_2=g\circ\Phi_1$ for some
 $g\in\Mob$, and $\langle1\rangle4$ applies; under (H),
 $\Sch(\Phi)$ is the corner measure by Theorem~\ref{thm:corner}.}
\lby{Lemma~\ref{lem:same-schwarz} and $\langle1\rangle4$}
\lqed{1}{6}
\end{proof}

\begin{remark}[Two caveats, both harmless for the state]\label{rem:caveat}
(i)  $\Phi$ is \emph{not} the zero-collar map of a single inclusion, so Box~[G1] is
applied per step only; normality of $\beta_\PP$ comes from
Theorem~\ref{thm:state}$\langle1\rangle2$, not from Thm.~7.1 applied to $\Phi$.
(ii) In a factorisation $\Phi=g\circ\Phi_1$ used in Theorem~\ref{thm:state}(2), the range of
$\beta_{V_{\Phi_1}}$ need not be contained in $\Fr(J_0)$ --- only $\Phi=g\circ\Phi_1$ is required
to map into $J_0$.  The \emph{state} $\omega\circ\beta_{V_{\Phi_1}}$ is nevertheless a well-defined
normal state on $\Fr(J_N)$ whenever $\Phi_1$ maps $J_N$ into some interval, which is all that
(2) asserts.  Both caveats were raised by REF-P6-0 (``Not checked'').
\end{remark}

\begin{corollary}[A single corner: the state depends only on $\zeta$, and not on the
side]\label{cor:phiA}
Let $I'=(\alpha,\omega_+)$ be a bounded interval, $p\in I'$, $\sigma>0$, and let
$k=\Rp_{p,\sigma}|_{I'}$ or $k=\Lp_{p,\sigma}|_{I'}$.  Then the state
$\omega\circ\beta_{V_k}$ on $\Fr(I')$, and hence
$\Phi_{\rm fid}=-\log F(\omega|_{\Fr(I')},\omega\circ\beta_{V_k})$, depends on
$(I',p,\sigma)$ and on the side of the compression \emph{only} through the invariant
$\zeta=\sigma\beta_{I'}(p)$ of Definition~\ref{def:strength}.  We write $\Phi_A(\zeta)$ for
this number (paper~1's single-step function, (C7)).
\end{corollary}
\begin{proof}
\lstep{1}{1}{Side-independence: $\Lp_{p,\sigma}=G\circ\Rp_{p,\sigma}$ with
 $G(x)=p+\frac{x-p}{1-\sigma(x-p)}\in\Mob$, so the two induce the same state.}
\lby{in the coordinate $w=1/(x-p)$ of Lemma~\ref{lem:parabolic-group}, $G$ is $w\mapsto w-\sigma$
 on both halves and $\Rp_{p,\sigma}$ is $w\mapsto w+\sigma$ on $\{w>0\}$, identity on
 $\{w<0\}$; the composite is the identity on $\{w>0\}$ and $w\mapsto w-\sigma$ on $\{w<0\}$,
 which is $\Lp_{p,\sigma}$; then Theorem~\ref{thm:state}(2)}
\lstep{1}{2}{$\Mob$ acts transitively on pairs (bounded interval, interior point): given
 $(I'_1,p_1)$ and $(I'_2,p_2)$ there is $g\in\Mob$, increasing and without pole on
 $\overline{I'_1}$, with $g(I'_1)=I'_2$, $g(p_1)=p_2$.}
\lby{$\Mob$ is $3$-transitive on $\mathbb R\cup\{\infty\}$ and preserves orientation and
 order: map the ordered triple (left endpoint, $p_1$, right endpoint) of $I'_1$ to that of
 $I'_2$; the pole then lies outside $\overline{I'_1}$ because the image of $\overline{I'_1}$ is
 the bounded set $\overline{I'_2}$}
\lstep{1}{3}{With such a $g$, $g\circ\Rp_{p_1,\sigma_1}\circ g^{-1}=\Rp_{p_2,\sigma_2}$ on
 $I'_2$, where $\sigma_2=\sigma_1/g'(p_1)$, and
 $\zeta_2=\sigma_2\beta_{I'_2}(p_2)=\sigma_1\beta_{I'_1}(p_1)=\zeta_1$.}
\lby{Lemma~\ref{lem:density}: conjugation by $g$ takes the parabolic flow at $p_1$ to the
 parabolic flow at $p_2$ with parameter $\sigma_1/g'(p_1)$, and maps the identity side to the
 identity side ($g$ increasing); and Lemma~\ref{lem:density}$\langle1\rangle5$ for the
 invariance of $\zeta$}
\lstep{1}{4}{Hence, writing $\Gamma:=\beta_{V_{g^{-1}}}:\Fr(I'_2)\to\Fr(I'_1)$, both
 $\omega|_{\Fr(I'_2)}$ and $\omega\circ\beta_{V_{k_2}}$ are the pullbacks along the
 $*$-isomorphism $\Gamma$ of $\omega|_{\Fr(I'_1)}$ and $\omega\circ\beta_{V_{k_1}}$.}
\lby{$\beta_{V_{k_2}}=\beta_{V_g}\beta_{V_{k_1}}\beta_{V_{g^{-1}}}$ by $\langle1\rangle3$ and
 Lemma~\ref{lem:comp-W}, together with $\omega\circ\beta_{V_g}=\omega$
 (Lemma~\ref{lem:moebius-inv})}
\lstep{1}{5}{$F$ is invariant under a common $*$-isomorphism, so the two fidelities agree and
 $\Phi_{\rm fid}$ is a function of $\zeta$ alone.}
\lby{Box~[G4], Def.~2: $P_{\rm tr}$ is defined from the $*$-algebra and its states only, hence
 is preserved by a $*$-isomorphism; $\langle1\rangle4$, $\langle1\rangle3$, $\langle1\rangle1$}
\lqed{1}{6}
\end{proof}

\section{Rigidity: exact recovery never happens (if the last step keeps a non-empty
part)}\label{sec:rigidity}

\begin{lemma}[Two-point rigidity, elementary]\label{lem:2pt}
Let $J=(\alpha,\beta)$ and let $\Phi:J\to\mathbb R$ be $C^1$ with $\Phi'>0$ and
\begin{equation}\label{eq:2pt}
 \frac{\sqrt{\Phi'(x)\Phi'(y)}}{\Phi(x)-\Phi(y)}=\frac1{x-y}\qquad\text{for all }x\ne y\in J .
\end{equation}
Then $\Phi$ is the restriction to $J$ of a M\"obius map without pole on $J$.  (Conversely every
such M\"obius map satisfies \eqref{eq:2pt}.)
\end{lemma}
\begin{proof}
\lstep{1}{1}{\eqref{eq:2pt} is equivalent to
 $(\Phi(x)-\Phi(y))^2=(x-y)^2\Phi'(x)\Phi'(y)$ for $x\ne y$.}
\lby{both sides of \eqref{eq:2pt} have the same sign ($\Phi$ increasing), so squaring and
 cross-multiplying is reversible}
\lstep{1}{2}{\lpick $y_1\in J$ and put $c:=\Phi'(y_1)>0$, $g:=\Phi-\Phi(y_1)$.  Then on
 $J_-:=(\alpha,y_1)$ and on $J_+:=(y_1,\beta)$: $g\ne0$, $g$ is $C^1$ and
 $g'(x)=\dfrac{g(x)^2}{c\,(x-y_1)^2}$.}
\lby{$\langle1\rangle1$ with $y=y_1$; $g\ne0$ because $\Phi$ is strictly increasing}
\lstep{1}{3}{On each of $J_\pm$, $\bigl(\tfrac1g\bigr)'=-\dfrac{g'}{g^2}
 =-\dfrac1{c(x-y_1)^2}=\Bigl(\dfrac1{c(x-y_1)}\Bigr)'$, hence there is a constant $A_\pm$ with
 $\dfrac1{g(x)}=\dfrac1{c(x-y_1)}+A_\pm$, i.e.
 $\Phi(x)=\Phi(y_1)+\dfrac{c(x-y_1)}{1+A_\pm c(x-y_1)}$.}
\lby{$\langle1\rangle2$; two functions with the same derivative on a connected interval differ by
 a constant}
\lstep{1}{4}{Hence $\Phi$ coincides on $J_-$ with a M\"obius map $M_-$ and on $J_+$ with a
 M\"obius map $M_+$, each without pole on the respective open interval.}
\lby{$\langle1\rangle3$: $x\mapsto\Phi(y_1)+c(x-y_1)/(1+A_\pm c(x-y_1))$ is M\"obius, and it is
 finite where $g$ is}
\lstep{1}{5}{$M_-=M_+$.}
\lproof
\lstep{2}{1}{\lpick $y_2\in J_+$.  Applying $\langle1\rangle3$--$\langle1\rangle4$ with $y_1$
 replaced by $y_2$ gives a M\"obius map $N_-$ with $\Phi=N_-$ on $(\alpha,y_2)$.}
\lby{$\langle1\rangle3$ used on the connected interval $(\alpha,y_2)$, which is one of the two
 components of $J\setminus\{y_2\}$}
\lstep{2}{2}{$N_-=M_-$ on $J_-$ and $N_-=M_+$ on $(y_1,y_2)$, two non-empty open intervals.}
\lby{$\langle2\rangle1$ and $\langle1\rangle4$; $(\alpha,y_2)\supseteq J_-\cup\{y_1\}\cup
 (y_1,y_2)$}
\lstep{2}{3}{Two M\"obius maps that agree on a non-empty open interval are equal.}
\lby{a M\"obius map is a rational function of degree $\le1$; the difference of two such, if zero
 on an open set, is zero as a rational function}
\lqed{2}{4} \lby{$\langle2\rangle2$, $\langle2\rangle3$: $M_-=N_-=M_+$}
\lstep{1}{6}{Therefore $\Phi=M_-=M_+$ on $J\setminus\{y_1\}$ and, by continuity, on $J$.}
\lby{$\langle1\rangle4$, $\langle1\rangle5$, $\Phi$ continuous}
\lstep{1}{7}{Converse: for $M$ M\"obius, $M(x)-M(y)=\frac{(ad-bc)(x-y)}{(cx+d)(cy+d)}$ and
 $M'(x)=\frac{ad-bc}{(cx+d)^2}$, so $(M(x)-M(y))^2=(x-y)^2M'(x)M'(y)$.}
\lby{direct computation}
\lqed{1}{8}
\end{proof}

\begin{lemma}[$F=1$ forces equality of states]\label{lem:F1}
Let $\varphi,\psi$ be states on a unital $C^*$- or $W^*$-algebra $\cM$ with
$P_{\rm tr}(\varphi,\psi)=1$.  Then $\varphi=\psi$.
\end{lemma}
\begin{proof}
\lstep{1}{1}{$1=P_{\rm tr}(\varphi,\psi)\le\varphi(a)\psi(a^{-1})$ for every invertible
 $a\in\cM_+$.}
\lby{Box~[G4], eq.~(1-4) of \cite{AU02}}
\lstep{1}{2}{\lpick $b=b^*\in\cM$ and $t\in\mathbb R$ with $|t|\,\lVert b\rVert<1/2$.  Then
 $a:=1+tb$ is positive and invertible and
 $\varphi(a)\psi(a^{-1})=1+t(\varphi(b)-\psi(b))+R(t)$ with $|R(t)|\le Ct^2$,
 $C=C(\lVert b\rVert)$.}
\lby{$a^{-1}=1-tb+t^2b^2(1+tb)^{-1}$ with $\lVert t^2b^2(1+tb)^{-1}\rVert\le2t^2\lVert
 b\rVert^2$, and $\varphi,\psi$ are states (norm one, unital)}
\lstep{1}{3}{$\langle1\rangle1$ and $\langle1\rangle2$ give
 $t(\varphi(b)-\psi(b))+R(t)\ge0$ for all small $|t|$, hence $\varphi(b)=\psi(b)$.}
\lby{dividing by $t>0$ and letting $t\downarrow0$, then by $t<0$ and letting $t\uparrow0$}
\lstep{1}{4}{Every element is a linear combination of two self-adjoint ones, so
 $\varphi=\psi$.}
\lby{$\langle1\rangle3$ and linearity}
\lqed{1}{5}
\end{proof}

\begin{theorem}[Exact recovery iff the map is M\"obius; and it never
happens]\label{thm:rigidity}
Let $\Phi\in\mathrm{PM}(I)$ with $\Phi(I)\subseteq\mathbb R$, and let
$\omega_\Phi:=\omega\circ\beta_{V_\Phi}$ be the induced state on $\Fr(I)$.  Then
\begin{enumerate}[label=(\arabic*),leftmargin=2.2em]
\item $\omega_\Phi=\omega|_{\Fr(I)}$ $\iff$ $\Phi$ is the restriction of a global M\"obius map
 $\iff$ $\Sch(\Phi)=0$;
\item for \emph{every} non-empty one-sided protocol $\PP$ whose last step keeps a non-empty part
 (C3) --- with or without (H) --- $\Sch(\Phi)\ne0$ on $I$; indeed $\Sch(\Phi)$ is a non-zero
 \emph{negative} measure, its total mass being at most $-2\sigma_N<0$;
\item hence $\omega_\PP:=\omega\circ\beta_\PP\ne\omega|_{\Fr(I)}$ and
 $\Phi_{\rm fid}(\PP)=-\log F(\omega|_{\Fr(I)},\omega_\PP)>0$ for every non-empty protocol
 \emph{whose last step keeps a non-empty part}.  The hypothesis cannot be dropped: see
 Remark~\ref{rem:degenerate} and REF-P6-3 \S2.
\end{enumerate}
\end{theorem}
\begin{proof}
\lstep{1}{1}{The two-point function of $\omega_\Phi$ has the distributional kernel
 $\sqrt{\Phi'(x)\Phi'(y)}\,Q(\Phi(x),\Phi(y))$, and that of $\omega$ is $Q(x,y)$, where
 $Q(x,y)=\frac1{2\pi i}\frac1{x-y-i0}\otimes\mathbf 1_r$ is the vacuum kernel.}
\lby{(C2), \eqref{eq:pushforward}: $\omega_\Phi(\psi(f)^*\psi(h))=\langle V_\Phi f,QV_\Phi
 h\rangle$; the Hardy kernel is the Cauchy kernel (Longo--Xu \cite{LX} \S3.1;
 \texttt{cited\_R1.tex} Box~C1)}
\lstep{1}{2}{If $\omega_\Phi=\omega$ then, testing with $f,h\in C^\infty_c(I)$ of \emph{disjoint}
 supports, $\sqrt{\Phi'(x)\Phi'(y)}/(\Phi(x)-\Phi(y))=1/(x-y)$ for all $x\ne y$ in $I$.}
\lby{$\langle1\rangle1$: off the diagonal both kernels are continuous functions, and two
 continuous functions defining the same distribution on the open set $\{x\ne y\}$ coincide there;
 the factor $\frac1{2\pi i}$ and the $i0$ are irrelevant off the diagonal, and $\Phi$ is
 increasing so $\Phi(x)-\Phi(y)$ has the sign of $x-y$}
\lstep{1}{3}{Hence $\Phi$ is M\"obius, by Lemma~\ref{lem:2pt}.}
\lby{$\langle1\rangle2$ and Lemma~\ref{lem:2pt} with $J=I$}
\lstep{1}{4}{Conversely, if $\Phi=g|_I$ with $g\in\Mob$ without pole on $I$, then
 $\omega_\Phi=\omega$ on $\Fr(I)$.}
\lby{Lemma~\ref{lem:moebius-inv}: $V_g$ commutes with $P$ and $V_gE_I=E_{g(I)}V_g$, so
 $E_IV_g^*PV_gE_I=E_IPE_I$}
\lstep{1}{5}{$\Sch(\Phi)=0\iff\Phi$ M\"obius.}
\lby{Lemma~\ref{lem:same-schwarz} with $\Phi_1=\id$ (whose Schwarzian vanishes): $\Sch(\Phi)=0$
 iff $\Phi=g\circ\id$; and Lemma~\ref{lem:masses} $\langle1\rangle1$ for the converse}
\lstep{1}{6}{(2): in \eqref{eq:general-rule} every term is a negative multiple of a point mass
 ($\sigma_m>0$, $G_m'>0$), so no cancellation is possible, and the term $m=N$ is always present
 because $G_N=\id$ and $x_N=p^{(N)}\in J_N$.}
\lby{Theorem~\ref{thm:general}; $G_N=\id$ by definition, and $p^{(N)}$ is an \emph{interior} point
 of $J_N=I$ because $A^{(N)}_B\subsetneq J_{N-1}$ (the last step keeps a non-empty part), so
 $p^{(N)}$ is the junction of $A^{(N)}_B$ with $J_{N-1}\setminus A^{(N)}_B\ne\emptyset$ and is
 therefore interior to $J_{N-1}\subseteq I$}
\lstep{1}{7}{(3): $\omega_\PP\ne\omega$ by (1) and (2); and $F=1$ would force
 $\omega_\PP=\omega$.}
\lby{Lemma~\ref{lem:F1} and $\Phi_{\rm fid}=-\log F$, $F\le1$}
\lqed{1}{8}
\end{proof}

\begin{remark}[The degenerate step, and why the hypothesis in (2) is needed]\label{rem:degenerate}
If a step conditions on the \emph{entire} current block ($A_B=J_{m-1}$, empty kept part), its
corner is the far endpoint of $J_{m-1}$.  For the last step this endpoint is an endpoint of
$I$, so the mass sits on $\partial I$ and $\Sch(\Phi)$ vanishes \emph{as a distribution on $I$}
if there is no other surviving corner; the extreme case $N=1$, $J_0=A_B$, is a pure M\"obius
compression of $I$ into $J_0$ and recovers the vacuum exactly.  This is consistent with
Definition~\ref{def:strength} and Lemma~\ref{lem:field}: $\bI(p)=0$ at $p\in\partial I$, so the
corner has strength $\zeta^{(I)}=0$; it is also Theorem~A at $a=0$, where
$\zeta=as/(a+L)=0$.  Such a step is not a recovery in the Petz sense --- it does not act as the
identity on any part of the block, so the ``recovered'' state does not agree with $\omega$ on any
conditioning algebra.
\end{remark}

\begin{verdictok}
This settles W1(iii) (``exact recovery never occurs'') \emph{unconditionally} and without the
second-order law: the brief's argument needed ``corner measure $\ne0\Rightarrow\Phi_{\rm fid}>0$''
(Theorem~A for one corner, part of W2 for several), whereas
Theorem~\ref{thm:rigidity} derives it from the M\"obius rigidity of the vacuum two-point function
plus the observation that all Schwarzian masses have the same sign.  Note that (2) also holds
when (H) fails and some corners are swallowed (Example~\ref{ex:swallow}): the \emph{last} step's
corner can never be swallowed.
\end{verdictok}

\section{The VWZ protocols, decided}\label{sec:vwz}

Throughout this section $n=4$ with $A=A_1$, $B=A_2$, $C=A_3$, $D=A_4$ of lengths $a,b,c,d$, and
$p_1<\dots<p_5$ as in (C1).  All statements are for an arbitrary common rotation parameter $t$;
$t=\infty$ ($s=\lambda$) and the ordinary Petz map $s=2\lambda$ ($t=0$) are the two ends of the
admissible family.

\begin{lemma}[The parabolic group law at a common corner]\label{lem:parabolic-group}
Fix $p$ and write $w:=1/(x-p)$.  In the coordinate $w$, $\Rp_{p,\sigma}$ is $w\mapsto w+\sigma$ on
$\{w>0\}$ and the identity on $\{w<0\}$; $\Lp_{p,\sigma}$ is $w\mapsto w-\sigma$ on $\{w<0\}$ and
the identity on $\{w>0\}$.  Consequently
\begin{align}
 \Rp_{p,\sigma_1}\circ\Rp_{p,\sigma_2}&=\Rp_{p,\sigma_1+\sigma_2},\qquad
 \Lp_{p,\sigma_1}\circ\Lp_{p,\sigma_2}=\Lp_{p,\sigma_1+\sigma_2},
 \label{eq:same-side}\\
 \Lp_{p,\sigma_1}\circ\Rp_{p,\sigma_2}&=\Rp_{p,\sigma_2}\circ\Lp_{p,\sigma_1}
 =M\circ\Rp_{p,\sigma_1+\sigma_2},\qquad M(x):=p+\frac{x-p}{1-\sigma_1(x-p)}\in\Mob .
 \label{eq:opposite-side}
\end{align}
\end{lemma}
\begin{proof}
\lstep{1}{1}{For $x>p$, $u:=x-p>0$: $\Rp_{p,\sigma}$ sends $u\mapsto u/(1+\sigma u)$, i.e.\
 $w=1/u\mapsto(1+\sigma u)/u=w+\sigma$, and $w+\sigma>0$; for $x<p$ it is the identity.}
\lby{\eqref{eq:step-map}}
\lstep{1}{2}{For $x<p$, $u<0$: $\Lp_{p,\sigma}$ sends $u\mapsto u/(1-\sigma u)$, i.e.\
 $w\mapsto w-\sigma$, and $w-\sigma<0$; for $x>p$ it is the identity.}
\lby{\eqref{eq:step-map}}
\lstep{1}{3}{\eqref{eq:same-side}: on $\{w>0\}$ the composite is $w\mapsto w+\sigma_1+\sigma_2$
 and on $\{w<0\}$ it is the identity.}
\lby{$\langle1\rangle1$; the half $\{w>0\}$ is preserved}
\lstep{1}{4}{$\Rp_{p,\sigma_2}$ and $\Lp_{p,\sigma_1}$ commute, and their composite is
 $w\mapsto w+\sigma_2$ on $\{w>0\}$, $w\mapsto w-\sigma_1$ on $\{w<0\}$.}
\lby{$\langle1\rangle1$, $\langle1\rangle2$: each map is the identity on the half the other moves,
 and each preserves its own half}
\lstep{1}{5}{$M$ is the global map $w\mapsto w-\sigma_1$ (both halves), and
 $M\circ\Rp_{p,\sigma_1+\sigma_2}$ is $w\mapsto w+\sigma_2$ on $\{w>0\}$, $w\mapsto w-\sigma_1$ on
 $\{w<0\}$; this is $\langle1\rangle4$.}
\lby{$1/(M(x)-p)=(1-\sigma_1(x-p))/(x-p)=w-\sigma_1$ for all $x\ne p$; $\langle1\rangle1$}
\lqed{1}{6}
\end{proof}

\begin{proposition}[VWZ Protocol 1(B) $=$ Protocol 1(A), literally]\label{prop:1B1A}
Protocol 1(A) is the single step $P_{B\to BCD}$: $A_B=B$, $A_{\rm new}=CD$, corner $p_2$
(here $A_{\rm new}$ is a union of two intervals of the fine chain, so 1(A) is a protocol of the
\emph{coarse} chain $A\mid B\mid CD$ in the sense of (C3); REF-P6-3 \S11),
parameter $\sigma_{1(A)}=\theta_t\,\frac{2(c+d)}{b(b+c+d)}$, where
$\theta_t:=\frac12(1+e^{-2\pi t})\in(\tfrac12,1]$ ($t\ge0$) and $\theta_0=1$ is the ordinary Petz value.  Protocol
1(B) is the two-step protocol $P_{B\to BC}$ then $P_{BC\to BCD}$, with corners
\[
 p^{(1)}=p_2,\quad \sigma_1=\theta_t\frac{2c}{b(b+c)},\qquad
 p^{(2)}=p_2,\quad \sigma_2=\theta_t\frac{2d}{(b+c)(b+c+d)} .
\]
Then $\sigma_1+\sigma_2=\sigma_{1(A)}$ and the two composite point maps are \emph{equal}:
$\Phi^{1(B)}=\Phi^{1(A)}=\Rp_{p_2,\sigma_{1(A)}}$.  In particular the recovered states agree, for
every common $t$, which is VWZ's assertion (Box~[G6]) with an exact proof.
\end{proposition}
\begin{proof}
\lstep{1}{1}{The corners and parameters are as stated.}
\lby{(C3): for $P_{B\to BC}$, $A_B=B$ and the corner is the end of $B$ away from $C$, i.e.\
 $p_2$; $\sigma=s/(l_B+l_{\rm new})$ with $s=2\theta_t\lambda$, $\lambda=c/b$, giving
 $\theta_t\frac{2c}{b(b+c)}$.  For $P_{BC\to BCD}$, $A_B=BC$ (a union, allowed by (C3)) and the
 corner is again $p_2$; $\lambda=d/(b+c)$, $l_B+l_{\rm new}=b+c+d$.  For $P_{B\to BCD}$,
 $\lambda=(c+d)/b$, $l_B+l_{\rm new}=b+c+d$}
\lstep{1}{2}{$\dfrac{2c}{b(b+c)}+\dfrac{2d}{(b+c)(b+c+d)}=\dfrac{2(c+d)}{b(b+c+d)}$.}
\lby{$\frac2{b+c}\Bigl[\frac cb+\frac d{b+c+d}\Bigr]
 =\frac2{b+c}\cdot\frac{c(b+c+d)+bd}{b(b+c+d)}$ and
 $c(b+c+d)+bd=c(b+c)+d(b+c)=(b+c)(c+d)$}
\lstep{1}{3}{Hence $\sigma_1+\sigma_2=\sigma_{1(A)}$ for every $t$ (the common factor $\theta_t$
 multiplies all three).}
\lby{$\langle1\rangle2$ and $\langle1\rangle1$}
\lstep{1}{4}{Both steps of 1(B) are right-compressions at the same point $p_2$, so
 $\Phi^{1(B)}=\Rp_{p_2,\sigma_1}\circ\Rp_{p_2,\sigma_2}=\Rp_{p_2,\sigma_1+\sigma_2}
 =\Rp_{p_2,\sigma_{1(A)}}=\Phi^{1(A)}$.}
\lby{\eqref{eq:same-side} and $\langle1\rangle3$}
\lqed{1}{5}
\end{proof}

\begin{remark}
Protocol 1(B) conditions on the \emph{union} $BC$, so it violates
Proposition~\ref{prop:H-suff}(a); it satisfies (H) through
Proposition~\ref{prop:H-suff}(b) (both corners at $p_2$).  This is the example that forced the
hypothesis of Theorem~\ref{thm:corner} to be stated in the weak form (H) (REF-P6-1).
\end{remark}

\begin{proposition}[VWZ Protocol 2]\label{prop:prot2}
Protocol 2 ($P_{B\to BC}$, then $P_{C\to CD}$) has corners $p_2$ with
$\sigma_2^{\rm VWZ}=\theta_t\frac{2c}{b(b+c)}$ and $p_3$ with
$\sigma_3^{\rm VWZ}=\theta_t\frac{2d}{c(c+d)}$, it satisfies
Proposition~\ref{prop:H-suff}(a), and
$\Sch(\Phi^{(2)})=-2\sigma_2^{\rm VWZ}\delta_{p_2}-2\sigma_3^{\rm VWZ}\delta_{p_3}$.
\end{proposition}
\begin{proof}
\lstep{1}{1}{Both conditioning regions are single chain intervals and $J_0=AB$ has two of them,
 so (H) holds.}
\lby{Proposition~\ref{prop:H-suff}(a)}
\lstep{1}{2}{Corners and parameters: as in Proposition~\ref{prop:1B1A}$\langle1\rangle1$ for the
 first step; for $P_{C\to CD}$, $A_B=C$, the corner is the end of $C$ away from $D$, i.e.\ $p_3$,
 $\lambda=d/c$ and $l_B+l_{\rm new}=c+d$.}
\lby{(C3)}
\lstep{1}{3}{The Schwarzian is as stated.}
\lby{Theorem~\ref{thm:corner} with $\langle1\rangle1$, $\langle1\rangle2$}
\lqed{1}{4}
\end{proof}

\begin{theorem}[VWZ Protocol 3 is an over-compressed single step]\label{thm:prot3}
Protocol 3 starts from $J_0=BC$ and performs $P_{B\to AB}$ (adjoining $A$, conditioning on $B$)
and $P_{C\to CD}$ (adjoining $D$, conditioning on $C$), in either order.  Then:
\begin{enumerate}[label=(\arabic*),leftmargin=2.2em]
\item both corners are at $p_3$, with
 $\sigma_1=\theta_t\frac{2a}{b(a+b)}$ (a left-compression) and
 $\sigma_2=\theta_t\frac{2d}{c(c+d)}$ (a right-compression);
\item the two orders give the \emph{same} point map, and
 $\Phi^{(3)}=M\circ\Rp_{p_3,\sigma_1+\sigma_2}$ with $M\in\Mob$ as in
 \eqref{eq:opposite-side}; hence $\Sch(\Phi^{(3)})=-2(\sigma_1+\sigma_2)\delta_{p_3}$;
\item the recovered state equals that of the single Theorem-A step with kept block $AB$,
 conditioning interval $C$, new interval $D$ and parameter
 $s'=(\sigma_1+\sigma_2)(c+d)$, which is admissible, indeed over-compressed:
 $s'>2\lambda'=2d/c$ \emph{strictly} (for $t=0$; $s'>2\theta_t\lambda'$ in general, and
 always $s'>\lambda'$ for $\theta_t\ge\frac12$, i.e.\ for every $t\ge0$);
\item therefore, exactly and for all sizes,
 \begin{equation}\label{eq:prot3}
  \Phi_{\rm fid}^{(3)}=\Phi_A(\zeta_{\rm eff}),\qquad
  \zeta_{\rm eff}=(\sigma_1+\sigma_2)\,\bI(p_3)
  =(\sigma_1+\sigma_2)\frac{(a+b)(c+d)}{a+b+c+d}.
 \end{equation}
\end{enumerate}
\end{theorem}
\begin{proof}
\lstep{1}{1}{(1): for $P_{B\to AB}$, $A_B=B$, $A_{\rm new}=A$ lies to the \emph{left}, so the
 corner is the right end of $B$, $p_3$, and $\sigma_1=s/(a+b)$ with $s=2\theta_t a/b$; for
 $P_{C\to CD}$ the corner is the left end of $C$, again $p_3$, with
 $\sigma_2=2\theta_t d/(c(c+d))$.}
\lby{(C3), \eqref{eq:corner-def}}
\lstep{1}{2}{(2): the two step maps are $\Lp_{p_3,\sigma_1}$ and $\Rp_{p_3,\sigma_2}$; they
 commute and their composite is $M\circ\Rp_{p_3,\sigma_1+\sigma_2}$.}
\lby{\eqref{eq:opposite-side}}
\lstep{1}{3}{$\Sch(\Phi^{(3)})=-2(\sigma_1+\sigma_2)\delta_{p_3}$.}
\lby{$\langle1\rangle2$, Lemma~\ref{lem:same-schwarz} ($M$ contributes nothing) and
 Lemma~\ref{lem:jump}; alternatively Theorem~\ref{thm:corner} with
 Proposition~\ref{prop:H-suff}(b)}
\lstep{1}{4}{(3): $\Rp_{p_3,\sigma_1+\sigma_2}$ is exactly the Theorem-A step with
 $A_{\rm TA}=AB$, $B_{\rm TA}=C$, $C_{\rm TA}=D$ and
 $\sigma=s'/(c+d)$, $s'=(\sigma_1+\sigma_2)(c+d)$; and
 $s'=\theta_t\bigl[\frac{2a(c+d)}{b(a+b)}+\frac{2d}c\bigr]>2\theta_t\frac dc$.}
\lby{(C3) with $l_B=c$, $l_{\rm new}=d$; and $\frac{2a(c+d)}{b(a+b)}>0$ for all
 $a,b,c,d>0$, whence the inequality is strict (REF-P6-3 \S10)}
\lstep{1}{5}{$s'\ge\lambda'=d/c$ for every $t\ge0$, so Box~[G1] applies to the
 effective single step and $\Phi_A$ is defined at $\zeta_{\rm eff}$.}
\lby{$\langle1\rangle4$ and $\theta_t=\frac12(1+e^{-2\pi t})\ge\frac12$ for $t\ge0$}
\lstep{1}{6}{(4): $\Phi^{(3)}$ and $\Rp_{p_3,\sigma_1+\sigma_2}$ induce the same state
 (Theorem~\ref{thm:state}(2) with $g=M$, which has its pole at $p_3+1/\sigma_1$, outside
 $\overline{\Rp(I)}$), and the strength of that single corner in the frame $I$ is
 $\zeta_{\rm eff}=(\sigma_1+\sigma_2)\bI(p_3)$ with
 $\bI(p_3)=\frac{(p_3-p_1)(p_5-p_3)}{p_5-p_1}=\frac{(a+b)(c+d)}{a+b+c+d}$.}
\lby{$\langle1\rangle2$, Theorem~\ref{thm:state}(2), Definition~\ref{def:strength},
 Lemma~\ref{lem:field} (which identifies $\zeta$ with paper~1's $as/(a+L)$: here
 $a=a+b$, $L=c+d$, $s=s'$); and Corollary~\ref{cor:phiA}, which is the step --- tacit in the
 first version of this document, REF-P6-3 \S10 --- that turns the identity of \emph{states} into
 the identity of \emph{numbers} $\Phi^{(3)}_{\rm fid}=\Phi_A(\zeta_{\rm eff})$, since for a
 single corner $\Phi_{\rm fid}$ is a function of $\zeta$ alone}
\lqed{1}{7}
\end{proof}

\begin{corollary}[Holonomy $=$ corner measure]\label{cor:holonomy}
Two one-sided protocols on the same chain that satisfy (H) and have the same corner measure
produce the same recovered state on $\Fr(I)$, hence the same fidelity, relative entropy and every
other state functional.  In particular Protocol 3 in either order, and Protocols 1(A) and 1(B),
are indistinguishable.
\end{corollary}
\begin{proof}
\lstep{1}{1}{Equal corner measures give equal Schwarzians.}
\lby{Theorem~\ref{thm:corner}}
\lstep{1}{2}{Equal Schwarzians give equal states.}
\lby{Theorem~\ref{thm:state}(3)}
\lqed{1}{3}
\end{proof}

\begin{verdictok}
Consequences (i) and (ii) of W1 are proved, with two strengthenings: 1(B) $=$ 1(A) holds for every
common rotation parameter $t$ and as an identity of \emph{point maps} (not only of states), and
the two orders of Protocol 3 give literally the same point map
(Lemma~\ref{lem:parabolic-group}), so no M\"obius quotient is needed for that half of the
statement.  The M\"obius quotient \emph{is} needed to compare Protocol 3 with the single
over-compressed step, exactly as REF-P6-0 F1b requires.  Numerics: \texttt{g1a\_calculus.py} \S4
(the identity in exact rational arithmetic; the maps agree to $2\cdot10^{-16}$; $s'>2\lambda'$
in $20\,000$ random rational configurations).
\end{verdictok}

\section{All one-sided protocols for $n=4$ and $n=5$}\label{sec:tables}

\begin{proposition}[With single-interval conditioning the holonomy is the starting
block]\label{prop:startblock}
Let $\PP$ be a complete one-sided protocol on $A_1,\dots,A_n$ in which every conditioning region
is a single chain interval, with starting block $J_0=A_i\cup\dots\cup A_j$, $j\ge i+1$, and a
common rotation parameter $t$, $\theta_t=\frac12(1+e^{-2\pi t})$.  Then (H) holds and
\begin{equation}\label{eq:startblock}
 \kappa_\PP=\sum_{k<i}\sigma^L_k\,\delta_{p_{k+2}}+\sum_{k>j}\sigma^R_k\,\delta_{p_{k-1}},
 \qquad
 \sigma^L_k=\theta_t\frac{2l_k}{l_{k+1}(l_k+l_{k+1})},\quad
 \sigma^R_k=\theta_t\frac{2l_k}{l_{k-1}(l_{k-1}+l_k)} ,
\end{equation}
\emph{independently of the order of the steps}.  Consequently all $\binom{(i-1)+(n-j)}{i-1}$
admissible orders for a given starting block give the same recovered state, and the number of distinct corner measures
among the $2^{n-2}$ protocols with a starting \emph{pair} is exactly $n-1$, one per pair.  The
left corners occupy $p_3,\dots,p_{i+1}$ and the right corners $p_j,\dots,p_{n-1}$; they overlap
in exactly one point, $p_{i+1}=p_j$, iff $j=i+1$, which then carries a double mass.
\end{proposition}
\begin{proof}
\lstep{1}{1}{(H) holds.}
\lby{Proposition~\ref{prop:H-suff}(a): single-interval conditioning and $|J_0|\ge2$ chain
 intervals}
\lstep{1}{2}{The step adjoining $A_k$ with $k>j$ has $A_B=A_{k-1}$, corner $p_{k-1}$ and
 $\sigma=\sigma^R_k$; the step adjoining $A_k$ with $k<i$ has $A_B=A_{k+1}$, corner $p_{k+2}$ and
 $\sigma=\sigma^L_k$.  Neither depends on the rest of the block, i.e.\ on the order.}
\lby{(C3): at the moment $A_k$ ($k>j$) is adjoined, the current block ends with $A_{k-1}$
 (only $A_{j+1},\dots,A_{k-1}$ can have been added on the right, and they are added in increasing
 order), so $A_B=A_{k-1}$ and the corner is the end of $A_{k-1}$ away from $A_k$, i.e.\
 $p_{k-1}$; $\lambda=l_k/l_{k-1}$, $l_B+l_{\rm new}=l_{k-1}+l_k$, $s=2\theta_t\lambda$.
 Symmetrically on the left}
\lstep{1}{3}{Hence the multiset $\{(p^{(m)},\sigma_m)\}$ is determined by $(i,j)$ and
 \eqref{eq:startblock} follows from Theorem~\ref{thm:corner}; the recovered state follows by
 Corollary~\ref{cor:holonomy}.}
\lby{$\langle1\rangle2$, Theorem~\ref{thm:corner}, Corollary~\ref{cor:holonomy}}
\lstep{1}{4}{The counting: an order is a shuffle of the $i-1$ left moves (which must be performed
 in the order $A_{i-1},A_{i-2},\dots$) with the $n-j$ right moves (in the order
 $A_{j+1},A_{j+2},\dots$), so there are $\binom{(i-1)+(n-j)}{i-1}$ of them; summing over the
 starting pairs $j=i+1$, $i=1,\dots,n-1$, gives $\sum_i\binom{n-2}{i-1}=2^{n-2}$.}
\lby{elementary counting; the left moves must be consecutive-inwards because $A_B$ is a single
 chain interval adjacent to the block}
\lstep{1}{5}{The corner ranges and the overlap are read off \eqref{eq:startblock}:
 $k\le i-1\Rightarrow p_{k+2}\le p_{i+1}$ and $k\ge j+1\Rightarrow p_{k-1}\ge p_j$.}
\lby{\eqref{eq:startblock}}
\lstep{1}{6}{Distinct starting pairs give distinct corner \emph{measures}, so the $n-1$ classes
 are really distinct.}
\lby{Lemma~\ref{lem:distinct} below --- the corner \emph{supports} do \emph{not} separate the
 classes, since for every middle pair ($2\le i\le n-2$) the support is all of
 $\{p_3,\dots,p_{n-1}\}$; this step was missing in the first version and is REF-P6-3 \S6}
\lqed{1}{7}
\end{proof}

\begin{lemma}[Distinct starting pairs, distinct measures]\label{lem:distinct}
In the situation of Proposition~\ref{prop:startblock} with starting pairs $(i,i+1)$ and
$(i',i'+1)$, $1\le i<i'\le n-1$, the two corner measures differ.
\end{lemma}
\begin{proof}
\lstep{1}{1}{By \eqref{eq:startblock} with $j=i+1$, the mass of the pair $(i,i+1)$ at $p_m$
 ($3\le m\le n-1$) is
 $\sigma^L_{m-2}\mathbf 1[m\le i+1]+\sigma^R_{m+1}\mathbf 1[m\ge i+1]$; in addition there is
 mass $\sigma^R_3$ at $p_2$ when $i=1$, and mass $\sigma^L_{n-2}$ at $p_n$ when $i=n-1$.}
\lby{\eqref{eq:startblock}: the left steps $k=1,\dots,i-1$ put $\sigma^L_k$ at $p_{k+2}$, i.e.\
 at $p_m$ with $m=k+2\le i+1$; the right steps $k=i+2,\dots,n$ put $\sigma^R_k$ at $p_{k-1}$,
 i.e.\ at $p_m$ with $m=k-1\ge i+1$}
\lstep{1}{2}{\lcase $2\le i$ and $i+1\le n-1$.  At $m=i+1$ the pair $(i,i+1)$ has mass
 $\sigma^L_{i-1}+\sigma^R_{i+2}$, while $(i',i'+1)$ has only $\sigma^L_{i-1}$ there, because
 $m=i+1\le i'+1$ but $m\ge i'+1$ fails.  Since
 $\sigma^R_{i+2}=\theta_t\,2l_{i+2}/(l_{i+1}(l_{i+1}+l_{i+2}))>0$, the measures differ.}
\lby{$\langle1\rangle1$, $i<i'$, and all lengths positive with $\theta_t>0$}
\lstep{1}{3}{\lcase $i=1$: the pair $(1,2)$ has mass $\sigma^R_3>0$ at $p_2$ and every
 $(i',i'+1)$ with $i'\ge2$ has none there.  \lcase $i'=n-1$ with $i\le n-2$: symmetrically at
 $p_n$.  These cases together with $\langle1\rangle2$ exhaust $1\le i<i'\le n-1$.}
\lby{$\langle1\rangle1$}
\lqed{1}{4}
\end{proof}

\begin{table}[h]\centering\small
\caption{$n=4$: all $2^{2}=4$ one-sided protocols with single-interval conditioning (ordinary
Petz, $\theta_t=1$), grouped by starting pair.  L $=$ add on the left, R $=$ add on the right.
$\sigma^R_k=2l_k/(l_{k-1}(l_{k-1}+l_k))$, $\sigma^L_k=2l_k/(l_{k+1}(l_k+l_{k+1}))$.  Last column:
$\zeta^{(I)}_p=\kappa_p\bI(p)$ for equal lengths $l=(1,1,1,1)$.}
\label{tab:n4}
\begin{tabular}{@{}llll@{}}
\toprule
start & orders & corner measure $\kappa_\PP$ & $\zeta^{(I)}$ (equal lengths)\\
\midrule
$A_1A_2$ & RR & $\sigma^R_3\delta_{p_2}+\sigma^R_4\delta_{p_3}$ & $0.75$ at $p_2$, $1$ at $p_3$\\
$A_2A_3$ & LR, RL & $(\sigma^L_1+\sigma^R_4)\,\delta_{p_3}$ & $2$ at $p_3$\\
$A_3A_4$ & LL & $\sigma^L_1\delta_{p_3}+\sigma^L_2\delta_{p_4}$ & $1$ at $p_3$, $0.75$ at $p_4$\\
\bottomrule
\end{tabular}
\end{table}

\begin{table}[h]\centering\small
\caption{$n=5$: all $2^{3}=8$ one-sided protocols with single-interval conditioning, grouped by
starting pair; $2$ classes of $3$ orders and $2$ classes of a single order.  Numerical column: ordinary Petz,
$l=(1.0,1.3,0.7,1.9,1.1)$, $T=6.0$ (\texttt{g1a\_calculus.py} \S5).}
\label{tab:n5}
\begin{tabular}{@{}llll@{}}
\toprule
start & orders & corner measure $\kappa_\PP$ & $\zeta^{(I)}$ at $(p_2,p_3,p_4,p_5)$\\
\midrule
$A_1A_2$ & RRR &
 $\sigma^R_3\delta_{p_2}+\sigma^R_4\delta_{p_3}+\sigma^R_5\delta_{p_4}$ &
 $(0.4487,\,2.9614,\,0.5789,\,-)$\\
$A_2A_3$ & LRR, RLR, RRL &
 $(\sigma^L_1+\sigma^R_4)\delta_{p_3}+\sigma^R_5\delta_{p_4}$ &
 $(-,\,3.9101,\,0.5789,\,-)$\\
$A_3A_4$ & LLR, LRL, RLL &
 $\sigma^L_1\delta_{p_3}+(\sigma^L_2+\sigma^R_5)\delta_{p_4}$ &
 $(-,\,0.9487,\,3.3647,\,-)$\\
$A_4A_5$ & LLL &
 $\sigma^L_1\delta_{p_3}+\sigma^L_2\delta_{p_4}+\sigma^L_3\delta_{p_5}$ &
 $(-,\,0.9487,\,2.7857,\,0.2546)$\\
\bottomrule
\end{tabular}
\end{table}

\begin{table}[h]\centering\small
\caption{$n=4$, union-conditioning variants (VWZ's own list), with the status of (H).  Corners
are the inner ends of the conditioning regions; $\sigma$'s for the ordinary Petz map.}
\label{tab:n4union}
\begin{tabular}{@{}lllll@{}}
\toprule
name & steps & corners & (H)? & corner measure\\
\midrule
1(A) & $P_{B\to BCD}$ & $p_2$ & yes (c) &
 $\frac{2(c+d)}{b(b+c+d)}\delta_{p_2}$\\
1(B) & $P_{B\to BC}$, $P_{BC\to BCD}$ & $p_2,p_2$ & yes (b) &
 $\bigl[\frac{2c}{b(b+c)}+\frac{2d}{(b+c)(b+c+d)}\bigr]\delta_{p_2}=$ 1(A)\\
2 & $P_{B\to BC}$, $P_{C\to CD}$ & $p_2,p_3$ & yes (a) &
 $\frac{2c}{b(b+c)}\delta_{p_2}+\frac{2d}{c(c+d)}\delta_{p_3}$\\
3 & $P_{B\to AB}$, $P_{C\to CD}$ & $p_3,p_3$ & yes (a),(b) &
 $\bigl[\frac{2a}{b(a+b)}+\frac{2d}{c(c+d)}\bigr]\delta_{p_3}$\\
--- & $P_{C\to CD}$, $P_{BC\to ABC}$ & $p_3,p_4$ & \textbf{no} &
 mass displaced, Example~\ref{ex:displaced}\\
\bottomrule
\end{tabular}
\end{table}

\begin{remark}[Which orders coincide, in words]\label{rem:coincide}
For $n=4$ there are three classes: $\{$1(A), 1(B)$\}$ (one corner at $p_2$), Protocol 2 ($p_2$ and
$p_3$), Protocol 3 in either order (a double corner at $p_3$), and the two mirror images of
1(A)/1(B) and of Protocol 2 obtained by reflecting the chain.  For $n=5$ there are four classes,
one per starting pair, and within a class \emph{every} interleaving of the left and right moves
gives the identical state.  The order therefore carries no information at all as long as (H)
holds and the conditioning regions are single intervals: the ``recovery holonomy'' of Note~3 is
the corner measure \eqref{eq:startblock}, i.e.\ the starting block.  Order dependence
\emph{does} appear when (H) fails (Examples~\ref{ex:displaced} and \ref{ex:swallow}), and that is
the only mechanism by which it can appear in the vacuum sector.
\end{remark}

\section{Amplification, and the left-to-right Petz chain in the modular frame}\label{sec:ampl}

\begin{lemma}[Amplification]\label{lem:ampl}
Let $p$ be a corner with parameter $\sigma$ and let $I_{\rm step}\subseteq I$ be two intervals
containing $p$.  Then
\begin{equation}\label{eq:ampl-mono}
 \zeta^{(I)}=\sigma\bI(p)\ \ge\ \sigma\beta_{I_{\rm step}}(p)=\zeta^{(I_{\rm step})},
 \qquad\text{with equality iff }I=I_{\rm step}.
\end{equation}
If moreover $I_{\rm step}=(p_1,f)$ with $f\le p_{n+1}$ --- the case of a step whose kept block
reaches the left end of the chain, $f$ being the far end of the moving part --- then, with
$a:=p-p_1$, $b:=f-p$, $c:=p_{n+1}-f$,
\begin{equation}\label{eq:ampl}
 \frac{\zeta^{(I)}}{\zeta^{(I_{\rm step})}}=\frac{\bI(p)}{\beta_{I_{\rm step}}(p)}
 =\frac{(a+b)(b+c)}{b(a+b+c)}=1+z',\qquad
 z'=\frac{ac}{b(a+b+c)},
\end{equation}
and $z'$ is simultaneously (i) the cross ratio of the four points $(p_1,p,f,p_{n+1})$ and (ii) the
Markov variable $z$ of paper~1 \eqref{eq:eta-z-ref} for the coarse triple (kept block $\mid$
$A_B\cup A_{\rm new}$ $\mid$ rest of the chain).  In particular $z'\ge0$, with $z'=0$ iff
$f=p_{n+1}$, i.e.\ iff the step is the last one.
\end{lemma}
\begin{proof}
\lstep{1}{1}{$\beta_{(\alpha,\omega)}(p)$ is strictly increasing in $\omega$ and strictly
 decreasing in $\alpha$; hence $I_{\rm step}\subseteq I$ implies \eqref{eq:ampl-mono}.}
\lby{$\partial_\omega\frac{(p-\alpha)(\omega-p)}{\omega-\alpha}
 =\frac{(p-\alpha)^2}{(\omega-\alpha)^2}>0$ and, by the reflection
 $x\mapsto-x$, $\partial_\alpha\beta<0$; monotonicity in each endpoint separately}
\lstep{1}{2}{With $p_1=0$, $p=a$, $f=a+b$, $p_{n+1}=a+b+c$:
 $\bI(p)=\frac{a(b+c)}{a+b+c}$ and $\beta_{I_{\rm step}}(p)=\frac{ab}{a+b}$, so the ratio is
 $\frac{(a+b)(b+c)}{b(a+b+c)}$.}
\lby{\eqref{eq:modframe} evaluated at the two intervals}
\lstep{1}{3}{$\frac{(a+b)(b+c)}{b(a+b+c)}=1+\frac{ac}{b(a+b+c)}$.}
\lby{$(a+b)(b+c)-b(a+b+c)=ab+ac+b^2+bc-ab-b^2-bc=ac$}
\lstep{1}{4}{$z'$ is the cross ratio of $(p_1,p,f,p_{n+1})$ and equals the paper~1 variable $z$
 of the coarse triple with lengths $(a,b,c)$.}
\lby{the cross ratio of the ordered quadruple $(p_1,p,f,p_{n+1})$ is
 $\frac{(p-p_1)(p_{n+1}-f)}{(f-p)(p_{n+1}-p_1)}=\frac{ac}{b(a+b+c)}$, and
 \texttt{paper1/sec\_setting.tex} eq.~(eta-z) defines $z=\frac{ac}{b(a+b+c)}$ for the triple with
 lengths $(a,b,c)$}
\lstep{1}{5}{$z'\ge0$ with equality iff $c=0$, i.e.\ $f=p_{n+1}$.}
\lby{$a,b>0$, $c\ge0$}
\lqed{1}{6}
\end{proof}

\noindent Here \eqref{eq:eta-z-ref} refers to
$\eta=\frac{b(a+b+c)}{(a+b)(b+c)}$, $z=\frac1\eta-1=\frac{ac}{b(a+b+c)}$
(\texttt{paper1/sec\_setting.tex}, l.~42--44), reproduced for the record:
\begin{equation}\label{eq:eta-z-ref}
 \eta=\frac{b(a+b+c)}{(a+b)(b+c)}\in(0,1),\qquad z=\frac1\eta-1=\frac{ac}{b(a+b+c)} .
\end{equation}

\begin{definition}[The L$\to$R Petz chain]\label{def:LR}
The \emph{left-to-right chain} is the protocol $J_0=A_1A_2$, then adjoin $A_{k+1}$ conditioning on
$A_k$ for $k=2,\dots,n-1$ (ordinary Petz, $\theta_t=1$).  By
Proposition~\ref{prop:startblock} it satisfies (H), its corners are $p_2,\dots,p_{n-1}$, all
simple, and
\begin{equation}\label{eq:LR-data}
 \kappa_k=\frac{2l_{k+1}}{l_k(l_k+l_{k+1})},\qquad
 \zeta_k:=\zeta^{(I)}_k=\kappa_k\frac{L_{k-1}(T-L_{k-1})}T,\qquad
 y_k=\log\frac{L_{k-1}}{T-L_{k-1}} .
\end{equation}
We write $\Delta_k:=y_{k+1}-y_k>0$ (the modular length of $A_k$ in the frame $I$),
$r_k=l_{k+1}/l_k$, $u_k=l_k/L_{k-1}$, $v_k=l_k/(T-L_k)$, and $\Delta_n:=+\infty$.
\end{definition}

\begin{lemma}[Exact corner data of the L$\to$R chain]\label{lem:LRdata}
For $2\le k\le n-1$:
\begin{align}
 e^{\Delta_k}&=(1+u_k)(1+v_k), &
 (e^{\Delta_k}-1)\,\zeta_k&=\frac{2r_k(1+v_k)}{1+r_k}\ \ge\ \frac{2r_k}{1+r_k},
 \label{eq:LR1}\\
 \zeta_k&=\frac{\sinh(\Delta_{k+1}/2)}
 {\sinh(\Delta_k/2)\,\sinh\bigl((\Delta_k+\Delta_{k+1})/2\bigr)}, &
 \frac1{\zeta_k}&=\sinh^2\!\frac{\Delta_k}2\coth\frac{\Delta_{k+1}}2+\frac12\sinh\Delta_k,
 \label{eq:LR2}
\end{align}
with the conventions $\coth(\infty)=1$ and, for $k=n-1$,
$\zeta_{n-1}=e^{-\Delta_{n-1}/2}/\sinh(\Delta_{n-1}/2)=2/(e^{\Delta_{n-1}}-1)$.  Moreover
\begin{equation}\label{eq:LR3}
 e^{\Delta_k}\le1+\frac2{\zeta_k}\qquad(2\le k\le n-1),
\end{equation}
\emph{with equality if and only if $k=n-1$}; the inequality is strict for $2\le k\le n-2$.
(Def.~\ref{def:LR} sets $\Delta_n=+\infty$, so at $k=n-1$ the ``limit'' $\Delta_{k+1}\to\infty$
is the actual value; REF-P6-3 \S4.)
\end{lemma}
\begin{proof}
\lstep{1}{1}{$e^{\Delta_k}=\frac{L_k}{T-L_k}\cdot\frac{T-L_{k-1}}{L_{k-1}}
 =\bigl(1+\frac{l_k}{L_{k-1}}\bigr)\bigl(1+\frac{l_k}{T-L_k}\bigr)=(1+u_k)(1+v_k)$.}
\lby{\eqref{eq:LR-data}, $L_k=L_{k-1}+l_k$ and $T-L_{k-1}=(T-L_k)+l_k$}
\lstep{1}{2}{$e^{\Delta_k}-1=u_k+v_k+u_kv_k=\dfrac{l_kT}{L_{k-1}(T-L_k)}$.}
\lby{$\langle1\rangle1$ and
 $\frac{l_k}{L_{k-1}}+\frac{l_k}{T-L_k}+\frac{l_k^2}{L_{k-1}(T-L_k)}
 =\frac{l_k[(T-L_k)+L_{k-1}+l_k]}{L_{k-1}(T-L_k)}$ with $L_{k-1}+l_k=L_k$}
\lstep{1}{3}{$\zeta_k=\dfrac{2r_k}{l_k(1+r_k)}\cdot\dfrac{L_{k-1}(T-L_{k-1})}T$, hence
 $(e^{\Delta_k}-1)\zeta_k=\dfrac{2r_k(T-L_{k-1})}{(1+r_k)(T-L_k)}=\dfrac{2r_k(1+v_k)}{1+r_k}$.}
\lby{$\kappa_k=\frac{2l_{k+1}}{l_k(l_k+l_{k+1})}=\frac{2r_k}{l_k(1+r_k)}$, $\langle1\rangle2$,
 and $\frac{T-L_{k-1}}{T-L_k}=1+v_k$}
\lstep{1}{4}{\eqref{eq:LR1} follows, the inequality because $v_k\ge0$ (equality iff $v_k=0$,
 i.e.\ $k=n-1$).}
\lby{$\langle1\rangle3$}
\lstep{1}{5}{Write $x_j:=(p_j-p_1)/T\in[0,1]$, so that $y_j=\log\frac{x_j}{1-x_j}$, and
 $\zeta_k=\dfrac{2\,x_k(1-x_k)(x_{k+2}-x_{k+1})}{(x_{k+1}-x_k)(x_{k+2}-x_k)}$.}
\lby{\eqref{eq:LR-data} with $l_k=T(x_{k+1}-x_k)$, $l_{k+1}=T(x_{k+2}-x_{k+1})$,
 $l_k+l_{k+1}=T(x_{k+2}-x_k)$, $L_{k-1}=Tx_k$, $T-L_{k-1}=T(1-x_k)$}
\lstep{1}{6}{For any $j<j'$: $x_{j'}-x_j
 =\dfrac{2\,e^{(y_j+y_{j'})/2}\sinh\bigl((y_{j'}-y_j)/2\bigr)}{(1+e^{y_j})(1+e^{y_{j'}})}$, and
 $x_j(1-x_j)=\dfrac{e^{y_j}}{(1+e^{y_j})^2}$.}
\lby{$x_j=\frac{e^{y_j}}{1+e^{y_j}}$, so
 $x_{j'}-x_j=\frac{e^{y_{j'}}-e^{y_j}}{(1+e^{y_j})(1+e^{y_{j'}})}$ and
 $e^{y_{j'}}-e^{y_j}=2e^{(y_j+y_{j'})/2}\sinh\frac{y_{j'}-y_j}2$}
\lstep{1}{7}{Substituting $\langle1\rangle6$ into $\langle1\rangle5$, every factor
 $(1+e^{y_j})$ and every exponential cancels and
 $\zeta_k=\dfrac{\sinh(\Delta_{k+1}/2)}{\sinh(\Delta_k/2)\sinh((\Delta_k+\Delta_{k+1})/2)}$.}
\lby{numerator $2\frac{e^{y_k}}{(1+e^{y_k})^2}\cdot
 \frac{2e^{(y_{k+1}+y_{k+2})/2}\sinh(\Delta_{k+1}/2)}{(1+e^{y_{k+1}})(1+e^{y_{k+2}})}$;
 denominator $\frac{2e^{(y_k+y_{k+1})/2}\sinh(\Delta_k/2)}{(1+e^{y_k})(1+e^{y_{k+1}})}\cdot
 \frac{2e^{(y_k+y_{k+2})/2}\sinh((\Delta_k+\Delta_{k+1})/2)}{(1+e^{y_k})(1+e^{y_{k+2}})}$; the
 ratio of the exponential prefactors is
 $e^{y_k+(y_{k+1}+y_{k+2})/2-(y_k+y_{k+1})/2-(y_k+y_{k+2})/2}=1$}
\lstep{1}{8}{The second identity in \eqref{eq:LR2} follows by expanding
 $\sinh\frac{\Delta_k+\Delta_{k+1}}2
 =\sinh\frac{\Delta_k}2\cosh\frac{\Delta_{k+1}}2+\cosh\frac{\Delta_k}2\sinh\frac{\Delta_{k+1}}2$
 and dividing numerator and denominator by $\sinh(\Delta_{k+1}/2)$.}
\lby{$\langle1\rangle7$ and the addition theorem}
\lstep{1}{9}{For $k=n-1$, $y_{n+1}=+\infty$, so $\Delta_n=\infty$, $\coth(\Delta_n/2)=1$ and
 $1/\zeta_{n-1}=\sinh^2\frac{\Delta_{n-1}}2+\frac12\sinh\Delta_{n-1}
 =\frac{e^{\Delta_{n-1}}-1}2$.}
\lby{$\langle1\rangle8$; $\sinh^2\frac u2=\frac{\cosh u-1}2$ and
 $\frac{\cosh u-1}2+\frac{\sinh u}2=\frac{e^u-1}2$}
\lstep{1}{10}{\eqref{eq:LR3}: $\coth(\Delta_{k+1}/2)\ge1$ gives
 $\frac1{\zeta_k}\ge\sinh^2\frac{\Delta_k}2+\frac12\sinh\Delta_k=\frac{e^{\Delta_k}-1}2$,
 i.e.\ $e^{\Delta_k}\le1+2/\zeta_k$; since $\sinh^2\frac{\Delta_k}2>0$, equality holds iff
 $\coth(\Delta_{k+1}/2)=1$, i.e.\ iff $\Delta_{k+1}=\infty$, i.e.\ iff $k=n-1$.}
\lby{$\langle1\rangle8$ for $2\le k\le n-2$ (there $\Delta_{k+1}<\infty$ and
 $\coth(\Delta_{k+1}/2)>1$, so the inequality is strict), $\langle1\rangle9$ for $k=n-1$ (there
 $1/\zeta_{n-1}=(e^{\Delta_{n-1}}-1)/2$ exactly)}
\lqed{1}{11}
\end{proof}

\begin{verdictok}
\eqref{eq:LR1} is item~6 of REF-P6-0 (re-derived here from scratch); \eqref{eq:LR2} and
\eqref{eq:LR3} are new.  \eqref{eq:LR2} says that the corner strengths of the L$\to$R chain are
functions of the \emph{modular lengths alone}: $\zeta_k=\zeta(\Delta_k,\Delta_{k+1})$.  This is
what makes the separation problem of \S\ref{sec:sep} a one-dimensional recursion.  Numerics:
\texttt{g1a\_separation.py} \S1, $50\,000$ random chains, maximum relative error
$3.5\cdot10^{-10}$ for \eqref{eq:LR1}--\eqref{eq:LR2} and $0$ violations of \eqref{eq:LR3}.
\end{verdictok}

\section{The separation constraint (W3), proved with the sharp constant}\label{sec:sep}

\begin{theorem}[Weak corners must be far apart]\label{thm:sep}
Consider the L$\to$R ordinary-Petz chain of Definition~\ref{def:LR} and let
$0<\eps<2$.  Fix $k\in\{2,\dots,n-1\}$ and assume
\begin{equation}\label{eq:allweak}
 \zeta_j\le\eps\qquad\text{for every }j\ \text{with}\ k\le j\le n-1 .
\end{equation}
Then
\begin{equation}\label{eq:sep}
 \boxed{\ e^{\Delta_k}\ \ge\ \frac2\eps\ }
\end{equation}
and, for the last corner, $e^{\Delta_{n-1}}\ge1+\frac2\eps$.  More precisely, with
$\gamma_j:=\coth(\Delta_j/2)-1=2/(e^{\Delta_j}-1)$ one has the exact recursion bound
$\gamma_j\le\zeta_j\bigl(1+\tfrac12\gamma_{j+1}\bigr)$ and hence
$\gamma_j\le\frac{2\eps}{2-\eps}$ for all $j\ge k$.
\end{theorem}
\begin{proof}
\lstep{1}{1}{Fix $j$ with $k\le j\le n-1$ and put $E:=e^{\Delta_j}\ge1$,
 $c_{j+1}:=\coth(\Delta_{j+1}/2)=1+\gamma_{j+1}\ge1$ (with $c_n=1$, $\gamma_n=0$).  Then
 \[ \frac4{\zeta_j}=c_{j+1}\Bigl(E-2+\frac1E\Bigr)+E-\frac1E. \]}
\lby{\eqref{eq:LR2}: $\frac1{\zeta_j}=\sinh^2\frac{\Delta_j}2\,c_{j+1}+\frac12\sinh\Delta_j$,
 together with $\sinh^2\frac u2=\frac{e^u-2+e^{-u}}4$ and
 $\frac12\sinh u=\frac{e^u-e^{-u}}4$}
\lstep{1}{2}{Hence $(c_{j+1}+1)E=\frac4{\zeta_j}+2c_{j+1}-\frac{c_{j+1}-1}E$.}
\lby{$\langle1\rangle1$, collecting the terms in $E$ and in $1/E$}
\lstep{1}{3}{$\dfrac{c_{j+1}-1}E\le c_{j+1}-1$, so
 $(c_{j+1}+1)E\ge\frac4{\zeta_j}+2c_{j+1}-(c_{j+1}-1)=\frac4{\zeta_j}+c_{j+1}+1$, i.e.
 \[ E\ \ge\ 1+\frac{4}{\zeta_j\,(2+\gamma_{j+1})} . \]}
\lby{$E\ge1$ (because $\Delta_j>0$) and $c_{j+1}\ge1$; $\langle1\rangle2$;
 $c_{j+1}+1=2+\gamma_{j+1}$}
\lstep{1}{4}{Therefore
 $\gamma_j=\dfrac2{E-1}\le\dfrac{\zeta_j(2+\gamma_{j+1})}2
 =\zeta_j\bigl(1+\tfrac12\gamma_{j+1}\bigr)$.}
\lby{$\langle1\rangle3$ and $t\mapsto2/t$ decreasing on $(0,\infty)$}
\lstep{1}{5}{\lassume \eqref{eq:allweak}.  \lprove $\gamma_j\le\frac{2\eps}{2-\eps}$ for
 $k\le j\le n-1$.}
\lproof
\lstep{2}{1}{Downward induction on $j$, starting from $\gamma_n=0\le\frac{2\eps}{2-\eps}$.}
\lby{$\Delta_n=+\infty$, so $\gamma_n=2/(e^{\infty}-1)=0$}
\lstep{2}{2}{\lassume $\gamma_{j+1}\le\frac{2\eps}{2-\eps}$.  Then
 $\gamma_j\le\eps\bigl(1+\frac12\cdot\frac{2\eps}{2-\eps}\bigr)
 =\eps\cdot\frac{2-\eps+\eps}{2-\eps}=\frac{2\eps}{2-\eps}$.}
\lby{$\langle1\rangle4$, $\zeta_j\le\eps$ by \eqref{eq:allweak}, and the induction hypothesis}
\lqed{2}{3} \lby{$\langle2\rangle1$, $\langle2\rangle2$ and induction}
\lstep{1}{6}{Combining $\langle1\rangle3$ at $j=k$ with $\langle1\rangle5$ at $j=k+1$ and
 $\zeta_k\le\eps$:
 \[ e^{\Delta_k}\ \ge\ 1+\frac4{\eps\bigl(2+\frac{2\eps}{2-\eps}\bigr)}
 =1+\frac{4(2-\eps)}{\eps\bigl(2(2-\eps)+2\eps\bigr)}=1+\frac{4(2-\eps)}{4\eps}
 =1+\frac{2-\eps}\eps=\frac2\eps . \]}
\lby{$\langle1\rangle3$ is decreasing in $\gamma_{k+1}$ and in $\zeta_k$; $\langle1\rangle5$;
 arithmetic}
\lstep{1}{7}{For $k=n-1$: $\gamma_n=0$, so $\langle1\rangle3$ gives
 $e^{\Delta_{n-1}}\ge1+\frac2{\zeta_{n-1}}\ge1+\frac2\eps$ (in fact with equality in the first
 inequality, by \eqref{eq:LR2} at $\Delta_n=\infty$).}
\lby{$\langle1\rangle3$ with $j=n-1$, $\gamma_n=0$; and Lemma~\ref{lem:LRdata}
 $\langle1\rangle9$}
\lqed{1}{8}
\end{proof}

\begin{corollary}[General admissible parameters]\label{cor:sep-admissible}
Let the chain be as in Definition~\ref{def:LR} but with arbitrary admissible parameters
$s_j\in[\lambda_j,2\lambda_j]$ at each step, and let $\zeta_j$ be the resulting strengths.  If
$\zeta_j\le\eps<1$ for all $j\ge k$, then $e^{\Delta_k}\ge1/\eps$.
\end{corollary}
\begin{proof}
\lstep{1}{1}{$\zeta_j=\frac{s_j}{2\lambda_j}\zeta_j^{\rm Petz}$ with
 $\frac{s_j}{2\lambda_j}\in[\frac12,1]$, hence
 $\zeta_j^{\rm Petz}\le2\zeta_j\le2\eps$.}
\lby{(C3): $\sigma_j=s_j/(l_j+l_{j+1})$ is proportional to $s_j$, and
 Definition~\ref{def:strength} is linear in $\sigma_j$; $s_j\ge\lambda_j$}
\lstep{1}{2}{Apply Theorem~\ref{thm:sep} with $\eps$ replaced by $2\eps<2$:
 $e^{\Delta_k}\ge\frac2{2\eps}=\frac1\eps$.}
\lby{$\langle1\rangle1$ and Theorem~\ref{thm:sep} (the modular positions $\Delta_j$ do not depend
 on the $s_j$)}
\lqed{1}{3}
\end{proof}

\begin{verdictok}
\textbf{W3 is proved, and the constant $2$ is sharp as $\eps\to0$.}  REF-P6-0's numerical
minimisation suggested $\min e^{\Delta}\approx2/\eps$ but, as REF-P6-1 noted, a minimisation only
bounds the minimum from \emph{above}; Theorem~\ref{thm:sep} supplies the matching lower bound
$e^{\Delta_k}\ge2/\eps$.  REF-P6-3 computed the \emph{exact} constrained minimum (a closed
backward recursion in the modular variables, itself a consequence of Lemma~\ref{lem:LRdata}):
\[
 \min e^{\Delta_k}=6.8134,\ 20.0499,\ 66.6817,\ 200.0050,\ 666.6682
 \quad\text{at }\eps=0.3,\,0.1,\,0.03,\,0.01,\,0.003,
\]
i.e.\ $(\eps/2)\min e^{\Delta_k}=1.0220,\ 1.0025,\ 1.00022,\ 1.000025,\ 1.000002$
(\texttt{ref\_p6\_3\_separation.out}~(b)).  So the bound is \emph{not} attained at fixed $\eps$
--- it is exceeded by $2.2\%$ at $\eps=0.3$ --- and the constant $2$ is sharp only in the limit
$\eps\to0$.  REF-P6-0's $\eps=0.3$ entry $6.75$ lies \emph{below} the exact minimum $6.8134$ and
is a float-cancellation artifact (a naive search over the \emph{lengths} destroys $T-L_k$); it is
not quoted here.  The remaining four entries $20.05$, $66.68$, $200.0$, $666.7$ agree with the
exact minima.  Independent checks (\texttt{g1a\_separation.py} \S\S2--3): $400\,000$ random
chains give $\min(\eps/2)e^{\Delta_k}=0.999999996$ (i.e.\ $1$ to floating-point accuracy); a
constrained minimisation reproduces $20.05$ and $200.01$ at $\eps=0.1,0.01$; the referee's optimal
length vector gives $(\eps/2)e^{\Delta_2}=1.000027$; and the recursion
$\gamma_j\le\zeta_j(1+\frac12\gamma_{j+1})$ and the fixed point $\gamma\le2\eps/(2-\eps)$ hold in
$200\,000$ random chains with $0$ violations.  Note that only the corners to the \emph{right} of
$p_k$ are used, and that the hypothesis cannot be weakened to the pair $\{k,k+1\}$: with only
$\zeta_k,\zeta_{k+1}\le\eps$ the separation can stay $O(1)$ (REF-P6-0: $e^{\Delta}=3.96$ at
$\eps=0.01$), the escape being paid for by a large $\zeta$ elsewhere.
\end{verdictok}

\begin{remark}[What the bound means for the second-order law]\label{rem:sep-meaning}
Combining \eqref{eq:sep} with the upper bound \eqref{eq:LR3} gives, for an all-weak L$\to$R chain,
$\frac2\eps\le e^{\Delta_k}\le1+\frac2{\zeta_k}$ (the right inequality an equality iff $k=n-1$):
the modular separation of consecutive corners is
pinned to within a factor $\eps^{-1}\zeta_k^{-1}$-ish window, and in particular
$\Delta_k\ge\log(2/\eps)\to\infty$ as $\eps\to0$.  Hence in the second-order network law W2
(track G1b) the cross terms between \emph{distinct} corners of an all-weak L$\to$R chain are
$O\bigl(\eps^2\sup_{\Delta\ge\log(2/\eps)}|\hat G(\Delta)|\bigr)$, i.e.\ they are controlled by
the decay of $\hat G$ --- which is open (card \texttt{corner-correlation-kernel}); and coincident
corners ($\Delta=0$) are \emph{not} excluded, since they belong to different protocols
(Protocol~3), not to a chain.  We state this only as the consequence of
Theorem~\ref{thm:sep}; the law itself, and the uniformity of its remainder, are W2 and are not
proved here.
\end{remark}

\section{The type-III sequential bound (W6)}\label{sec:w6}

Nothing in this section uses the free fermion or the geometry: it holds for any chain of von
Neumann inclusions and any normal recovery channels.

\begin{hypothesis}[(U): Uhlmann's theorem in the standard form]\label{hyp:U}
Let $\cM$ be a $\sigma$-finite von Neumann algebra in standard form
$(\mathcal H,J,\mathcal P^\natural)$.  Every normal state $\varphi$ has a unique vector
representative $\xi_\varphi\in\mathcal P^\natural$, and for all normal states $\varphi,\psi$
\begin{equation}\label{eq:U}
 F(\varphi,\psi)=\sqrt{P_{\rm tr}(\varphi,\psi)}=\langle\xi_\varphi,\xi_\psi\rangle\in[0,1].
\end{equation}
\emph{Status} (corrected after REF-P6-3 \S8).  \eqref{eq:U} is the standard-form version of
Uhlmann's theorem, usually attributed to Araki--Raggio \cite{AR82}, which is \textbf{NOT
OBTAINED} (not in \texttt{refs/}).  Uhlmann \cite{Uhl76} \emph{is} on disk
(\texttt{refs/Uhlmann76.pdf}) --- an earlier version of this document wrongly called it not
obtained --- and was read: it contains the explicit-formula theorem (unnumbered, p.~277),
one-argument concavity of $P_{\rm tr}$ (eq.~(8), p.~274) and monotonicity under
\emph{restriction to a subalgebra} (eq.~(24), p.~277), but \emph{no} standard form, no natural
cone and no metric statement; Alberti--Uhlmann \cite{AU02} has no natural-cone formula either.
We therefore declare \eqref{eq:U} as a hypothesis rather than citing it as verified.  It is used
\emph{only} for the angle inequality \eqref{eq:w6} and the exact form \eqref{eq:w6-master}: the
second-order bound \eqref{eq:w6-2nd}, which is the form W6 needs, is proved without (U) in
Theorem~\ref{thm:w6chord} and Corollary~\ref{cor:w6-second}(2).  Only two consequences are
used: all the states occurring below have vector representatives \emph{in one and the same}
Hilbert space, and for those representatives the fidelity is their (non-negative, real) inner
product.  The upper bound $F\le1$ and the variational formula of Box~[G4] are \emph{not}
part of (U) and are used independently.
\end{hypothesis}

\begin{lemma}[Choi's inequality for unital $2$-positive maps]\label{lem:choi}
Let $T:\cN\to\cM$ be unital and $2$-positive and let $a\in\cN$ be positive and invertible.  Then
$T(a)$ is positive and invertible and $T(a^{-1})\ge T(a)^{-1}$.
\end{lemma}
\begin{proof}
\lstep{1}{1}{$a\ge\lVert a^{-1}\rVert^{-1}1$, so $T(a)\ge\lVert a^{-1}\rVert^{-1}1>0$ and $T(a)$
 is invertible.}
\lby{positivity and unitality of $T$}
\lstep{1}{2}{$\begin{pmatrix}a&1\\1&a^{-1}\end{pmatrix}
 =\begin{pmatrix}a^{1/2}\\a^{-1/2}\end{pmatrix}
 \begin{pmatrix}a^{1/2}&a^{-1/2}\end{pmatrix}\ge0$ in $M_2(\cN)$.}
\lby{it is of the form $v^*v$}
\lstep{1}{3}{Applying $T\otimes\id_{M_2}$:
 $\begin{pmatrix}T(a)&1\\1&T(a^{-1})\end{pmatrix}\ge0$.}
\lby{$T$ is $2$-positive and unital, so $T(1)=1$}
\lstep{1}{4}{For a positive block matrix $\begin{pmatrix}X&1\\1&Y\end{pmatrix}$ with $X>0$
 invertible, $Y\ge X^{-1}$.}
\lby{testing on the vector $(-X^{-1}\eta,\eta)$ gives
 $\langle\eta,X^{-1}\eta\rangle-2\langle\eta,X^{-1}\eta\rangle+\langle\eta,Y\eta\rangle\ge0$}
\lqed{1}{5}
\end{proof}

\begin{lemma}[Monotonicity of the fidelity]\label{lem:monotone}
Let $T:\cN\to\cM$ be a normal unital $2$-positive map and $\varphi,\psi$ normal states on $\cM$.
Then $F(\varphi\circ T,\psi\circ T)\ge F(\varphi,\psi)$, i.e.\ $\Bang$ decreases under channels.
\end{lemma}
\begin{proof}
\lstep{1}{1}{For every positive invertible $a\in\cN$:
 $(\varphi\circ T)(a)\,(\psi\circ T)(a^{-1})=\varphi(T(a))\,\psi(T(a^{-1}))
 \ge\varphi(T(a))\,\psi(T(a)^{-1})\ge P_{\rm tr}(\varphi,\psi)$.}
\lby{Lemma~\ref{lem:choi} and positivity of $\psi$ for the first inequality; Box~[G4]
 eq.~(1-4) with the positive invertible element $T(a)\in\cM$ for the second}
\lstep{1}{2}{Taking the infimum over $a$ and using the variational formula
 $P_{\rm tr}=\inf_a\varphi(a)\psi(a^{-1})$ on $\cN$ gives
 $P_{\rm tr}(\varphi\circ T,\psi\circ T)\ge P_{\rm tr}(\varphi,\psi)$.}
\lby{$\langle1\rangle1$ and Box~[G4] (Alberti--Uhlmann Thm.~2(1))}
\lqed{1}{3}
\end{proof}

\begin{lemma}[Spherical triangle inequality]\label{lem:spherical}
Let $x,y,z$ be unit vectors in a real or complex Hilbert space with $\langle x,y\rangle$,
$\langle y,z\rangle$, $\langle x,z\rangle$ real and non-negative, and put
$\alpha=\arccos\langle x,y\rangle$, $\beta=\arccos\langle y,z\rangle$,
$\gamma=\arccos\langle x,z\rangle\in[0,\pi/2]$.  Then $\gamma\le\alpha+\beta$.
\end{lemma}
\begin{proof}
\lstep{1}{1}{\lcase $\alpha+\beta\ge\pi/2$: then $\gamma\le\pi/2\le\alpha+\beta$.}
\lby{$\langle x,z\rangle\ge0$ gives $\gamma\le\pi/2$}
\lstep{1}{2}{\lcase $\alpha+\beta<\pi/2$.  Write $x=\cos\alpha\,y+\sin\alpha\,u$ and
 $z=\cos\beta\,y+\sin\beta\,v$ with $u,v$ unit vectors orthogonal to $y$ (if $\sin\alpha=0$ then
 $x=y$ and the claim is trivial; likewise for $\beta$).}
\lby{orthogonal decomposition along $y$: $x=\langle y,x\rangle y+x^\perp$ with
 $\lVert x^\perp\rVert=\sqrt{1-\cos^2\alpha}=\sin\alpha\ge0$}
\lstep{1}{3}{$\langle x,z\rangle=\cos\alpha\cos\beta+\sin\alpha\sin\beta\,
 \Re\langle u,v\rangle\ge\cos\alpha\cos\beta-\sin\alpha\sin\beta=\cos(\alpha+\beta)$.}
\lby{$\langle1\rangle2$, Cauchy--Schwarz $|\langle u,v\rangle|\le1$, and
 $\sin\alpha,\sin\beta\ge0$}
\lstep{1}{4}{$\arccos$ is decreasing, so $\gamma\le\alpha+\beta$.}
\lby{$\langle1\rangle3$ and $\alpha+\beta\in[0,\pi]$}
\lqed{1}{5}
\end{proof}

\begin{theorem}[Sequential recovery along a chain of von Neumann inclusions]\label{thm:w6}
Assume (U).  Let $\cM$ be a von Neumann algebra with a normal state $\omega$, let
$\cN_1\subseteq\cN_2\subseteq\dots\subseteq\cN_m\subseteq\cM$ be von Neumann subalgebras with the
same unit, and let $R_k:\cN_{k+1}\to\cN_k$ ($k=1,\dots,m-1$) be normal unital $2$-positive
(Heisenberg-picture) recovery maps.  Put $R:=R_1\circ R_2\circ\dots\circ R_{m-1}:\cN_m\to\cN_1$.
Then, with $\Bang(\varphi,\psi)=\arccos F(\varphi,\psi)$,
\begin{equation}\label{eq:w6}
 \boxed{\ \Bang\bigl(\omega|_{\cN_m},\ \omega|_{\cN_1}\circ R\bigr)\ \le\
 \sum_{k=1}^{m-1}\Bang\bigl(\omega|_{\cN_{k+1}},\ \omega|_{\cN_k}\circ R_k\bigr).\ }
\end{equation}
\end{theorem}
\begin{proof}
\lstep{1}{1}{\ldefine $S_j:=R_{m-j+1}\circ\dots\circ R_{m-1}:\cN_m\to\cN_{m-j+1}$ for
 $j=1,\dots,m-1$ ($S_1=\id_{\cN_m}$), and the normal states on $\cN_m$
 \[ \mu_0:=\omega|_{\cN_m},\qquad
 \mu_j:=\omega|_{\cN_{m-j}}\circ R_{m-j}\circ S_j\quad(j=1,\dots,m-1). \]}
\lby{each $R_k$ is normal unital positive, so each composite maps $\cN_m$ into $\cN_{m-j}$ and
 each $\mu_j$ is a normal state}
\lstep{1}{2}{$\mu_{m-1}=\omega|_{\cN_1}\circ R$ and $\mu_{j-1}=\omega|_{\cN_{m-j+1}}\circ S_j$.}
\lby{$S_{m-1}=R_2\circ\dots\circ R_{m-1}$, so
 $\mu_{m-1}=\omega|_{\cN_1}\circ R_1\circ S_{m-1}=\omega|_{\cN_1}\circ R$; and
 $\mu_{j-1}=\omega|_{\cN_{m-j+1}}\circ R_{m-j+1}\circ S_{j-1}=\omega|_{\cN_{m-j+1}}\circ S_j$ for
 $j\ge2$, while for $j=1$ it is $\mu_0$ with $S_1=\id$}
\lstep{1}{3}{$\Bang(\mu_{j-1},\mu_j)\le
 \Bang\bigl(\omega|_{\cN_{m-j+1}},\omega|_{\cN_{m-j}}\circ R_{m-j}\bigr)$.}
\lby{$\langle1\rangle2$: both states are the images of
 $\omega|_{\cN_{m-j+1}}$ and $\omega|_{\cN_{m-j}}\circ R_{m-j}$ (two normal states on
 $\cN_{m-j+1}$) under the same channel $S_j$, and Lemma~\ref{lem:monotone} ($F$ increases, hence
 $\Bang$ decreases)}
\lstep{1}{4}{$\Bang(\mu_0,\mu_{m-1})\le\sum_{j=1}^{m-1}\Bang(\mu_{j-1},\mu_j)$.}
\lby{Hypothesis~\ref{hyp:U}: all $\mu_j$ have natural-cone representatives $\xi_j$ in the same
 standard form of $\cN_m$, $F(\mu_i,\mu_j)=\langle\xi_i,\xi_j\rangle\ge0$; hence
 $\Bang$ is the angle between unit vectors and Lemma~\ref{lem:spherical} applies, by induction on
 $m$}
\lstep{1}{5}{Combining $\langle1\rangle4$ with $\langle1\rangle3$ and re-indexing $k=m-j$ gives
 \eqref{eq:w6}.}
\lby{$\langle1\rangle2$--$\langle1\rangle4$}
\lqed{1}{6}
\end{proof}

\begin{citedbox}{[G7] Alberti--Uhlmann 2002 Prop.~1(1), p.~3: the Bures distance is a metric}
\emph{Transcribed verbatim} (\texttt{refs/AlbertiUhlmann02\_math-ph-0202038.pdf}, p.~3):
``\textsc{Proposition 1}.  Let $d_{\rm B}:M^*_+\times M^*_+\ni\{\nu,\varrho\}\mapsto
d_{\rm B}(M|\nu,\varrho)\in\mathbb R_+$ be given in accordance with Definition 1.  Then the
following hold: (1) $d_{\rm B}$ is a distance function on the points of $M^*_+$; (2)
$d_{\rm B}$ is topologically equivalent with $d_1$ on bounded subsets of $M^*_+$.''  Definition~1,
p.~1, is $d_{\rm B}(M|\nu,\varrho)=\inf\lVert\psi-\phi\rVert$ over representations and vector
representatives, and eq.~(1-3), p.~4, gives
$d_{\rm B}^2=(\lVert\nu\rVert_1-\sqrt{P_M})+(\lVert\varrho\rVert_1-\sqrt{P_M})$, i.e.\
$d_{\rm B}=\sqrt{2-2F}$ for states.

\emph{Caveats.}  The proposition is stated \emph{without proof} (``The relevant basic facts will
be stated here without proof'') and item~(1) is attributed to D.~Bures, Trans.\ Amer.\ Math.\
Soc.\ \textbf{135} (1969) 199--212, which is \textbf{NOT OBTAINED} (not in \texttt{refs/}).  The
metric is the \emph{chordal} one; neither \cite{AU02} nor \cite{Uhl76} says anything about the
Bures \emph{angle}.

\emph{Use here.}  Theorem~\ref{thm:w6chord}: it replaces Hypothesis~\ref{hyp:U} in the only
statement W6 needs.

\smallskip\noindent\shot[\textwidth]{G1a_AU02_prop1_p3.png}
\end{citedbox}

\begin{theorem}[Chordal sequential bound; no Hypothesis (U)]\label{thm:w6chord}
In the setting of Theorem~\ref{thm:w6} but \emph{without assuming} Hypothesis~(U) (only the standing
hypotheses on the algebras, the normal state $\omega$ and the channels $R_k$), with
$d_{\rm B}(\varphi,\psi):=\sqrt{2-2F(\varphi,\psi)}$ the Bures distance on normal states,
\begin{equation}\label{eq:w6chord}
 d_{\rm B}\bigl(\omega|_{\cN_m},\ \omega|_{\cN_1}\circ R\bigr)\ \le\
 \sum_{k=1}^{m-1}d_{\rm B}\bigl(\omega|_{\cN_{k+1}},\ \omega|_{\cN_k}\circ R_k\bigr).
\end{equation}
\end{theorem}
\begin{proof}
\noindent Steps $\langle1\rangle1$--$\langle1\rangle3$ of Theorem~\ref{thm:w6} do not use Hypothesis~(U);
only its step $\langle1\rangle4$ does, and that step is not invoked here.
\lstep{1}{1}{$d_{\rm B}$ satisfies the triangle inequality on normal states, and
 $d_{\rm B}=\sqrt{2-2F}$ there.}
\lby{Box~[G7]: \cite{AU02} Prop.~1(1), p.~3, and eq.~(1-3), p.~4, with
 $\lVert\nu\rVert_1=\lVert\varrho\rVert_1=1$}
\lstep{1}{2}{With the states $\mu_0,\dots,\mu_{m-1}$ of Theorem~\ref{thm:w6}
 $\langle1\rangle1$, $d_{\rm B}(\mu_0,\mu_{m-1})\le\sum_{j}d_{\rm B}(\mu_{j-1},\mu_j)$.}
\lby{$\langle1\rangle1$ and induction on $m$; all $\mu_j$ are normal states on $\cN_m$}
\lstep{1}{3}{$d_{\rm B}(\mu_{j-1},\mu_j)\le
 d_{\rm B}(\omega|_{\cN_{m-j+1}},\omega|_{\cN_{m-j}}\circ R_{m-j})$.}
\lby{Theorem~\ref{thm:w6} $\langle1\rangle2$ (both states are images under the same channel
 $S_j$), Lemma~\ref{lem:monotone} ($F$ increases under $S_j$) and the fact that
 $t\mapsto\sqrt{2-2t}$ is decreasing}
\lstep{1}{4}{Re-indexing $k=m-j$ gives \eqref{eq:w6chord}.}
\lby{$\langle1\rangle2$, $\langle1\rangle3$}
\lqed{1}{5}
\end{proof}

\begin{corollary}[Second-order form]\label{cor:w6-second}
In the setting of Theorem~\ref{thm:w6}, Hypothesis~(U) being assumed only where stated, put $\Phi_{\rm tot}:=-\log F(\omega|_{\cN_m},
\omega|_{\cN_1}\circ R)$ and $\Phi_k:=-\log F(\omega|_{\cN_{k+1}},\omega|_{\cN_k}\circ R_k)$.
Then: \emph{(1) under Hypothesis~\ref{hyp:U}},
\begin{equation}\label{eq:w6-master}
 \Phi_{\rm tot}\le-\log\cos\Bigl(\sum_k\arccos e^{-\Phi_k}\Bigr)
 \quad\text{whenever }\sum_k\arccos e^{-\Phi_k}<\frac\pi2,
\end{equation}
and \emph{(2) unconditionally} (no Hypothesis~(U)), as all $\Phi_k\to0$,
\begin{equation}\label{eq:w6-2nd}
 \Phi_{\rm tot}\le\Bigl(\sum_k\sqrt{\Phi_k}\Bigr)^2\bigl(1+o(1)\bigr).
\end{equation}
If, in addition, each step map extends to a normal unital $2$-positive map
$\check R_k:\Fr(I)\to\Fr(I)$ with $\check R_k|_{\cN_{k+1}}=R_k$ and error
$\check\Phi_k:=-\log F(\omega,\omega\circ\check R_k)$ measured on all of $\Fr(I)$, then
$\Phi_k\le\check\Phi_k$ and \eqref{eq:w6-2nd} holds with $\check\Phi_k$ in place of $\Phi_k$.
\end{corollary}
\begin{proof}
\lstep{1}{1}{$\Bang=\arccos e^{-\Phi_{\rm fid}}$ and $\Phi_{\rm fid}=-\log\cos\Bang$ are
 increasing bijections between $[0,\infty)$ and $[0,\pi/2)$.}
\lby{(C7): $F=e^{-\Phi_{\rm fid}}=\cos\Bang$}
\lstep{1}{2}{\eqref{eq:w6-master} follows from \eqref{eq:w6} by applying the increasing function
 $\Bang\mapsto-\log\cos\Bang$ on $[0,\pi/2)$.}
\lby{$\langle1\rangle1$, Theorem~\ref{thm:w6}, and monotonicity}
\lstep{1}{3}{As $\Phi\downarrow0$, $\arccos e^{-\Phi}=\sqrt{2\Phi}\,(1+O(\Phi))$ and
 $-\log\cos\Bang=\frac12\Bang^2(1+O(\Bang^2))$; substituting into \eqref{eq:w6-master} gives
 \eqref{eq:w6-2nd}, under (U).}
\lby{$e^{-\Phi}=1-\Phi+O(\Phi^2)$ and $\arccos(1-t)=\sqrt{2t}(1+O(t))$;
 $-\log\cos A=\frac{A^2}2+O(A^4)$; and
 $\frac12\bigl(\sum_k\sqrt{2\Phi_k}\bigr)^2=\bigl(\sum_k\sqrt{\Phi_k}\bigr)^2$}
\lstep{1}{3a}{\eqref{eq:w6-2nd} again, \emph{without} (U): by Theorem~\ref{thm:w6chord},
 $\sqrt{2-2e^{-\Phi_{\rm tot}}}\le\sum_k\sqrt{2-2e^{-\Phi_k}}$, and
 $\sqrt{2-2e^{-\Phi}}=\sqrt{2\Phi}\,(1+O(\Phi))$ on both sides; squaring and dividing by $2$
 gives $\Phi_{\rm tot}\le\bigl(\sum_k\sqrt{\Phi_k}\bigr)^2(1+o(1))$.}
\lby{Theorem~\ref{thm:w6chord}, $F=e^{-\Phi_{\rm fid}}$ (C7), $2-2e^{-\Phi}=2\Phi+O(\Phi^2)$,
 and $t\mapsto-\log(1-t^2/2)$ increasing; numerically the relative excess of the exact bound over
 $(\sum\sqrt{\Phi_k})^2$ falls as $1.7\cdot10^{-1}\to9.6\cdot10^{-4}\to1.0\cdot10^{-5}$ for
 $\Phi_k\le10^{-2},10^{-4},10^{-6}$ (\texttt{ref\_p6\_3\_separation.out}~(d)), against
 $1.1\cdot10^{-1}\to7.2\cdot10^{-4}\to7.4\cdot10^{-6}$ for the angle route}
\lstep{1}{4}{$\Phi_k\le\check\Phi_k$ under the extension hypothesis.}
\lby{restriction to the subalgebra $\cN_{k+1}\subseteq\Fr(I)$ is a normal unital
 $*$-homomorphism, hence a channel, so Lemma~\ref{lem:monotone} gives
 $F(\omega|_{\cN_{k+1}},\omega|_{\cN_k}\circ R_k)\ge F(\omega,\omega\circ\check R_k)$}
\lqed{1}{5}
\end{proof}

\begin{remark}[Saturation, and the comparison with the network law]\label{rem:w6-sat}
For a protocol all of whose corners coincide at one point $p$ (e.g.\ VWZ 1(B) or Protocol 3),
Theorem~\ref{thm:prot3} and Theorem~A give $\Phi_{\rm tot}=\Phi_A\bigl(\sum_k\zeta_k\bigr)
=f_2\bigl(\sum_k\zeta_k\bigr)^2+o$, while each step contributes $\check\Phi_k=f_2\zeta_k^2+o$;
so \eqref{eq:w6-2nd} is \emph{saturated at second order}.  For corners at distinct points the
bound is far from sharp: the conjectural second-order network law W2 (track G1b) replaces
$\bigl(\sum\zeta_k\bigr)^2$ by $\sum_{j,k}\zeta_j\zeta_k\hat G(y_j-y_k)/\hat G(0)$, and
Theorem~\ref{thm:sep} shows that for an all-weak L$\to$R chain the off-diagonal terms are
evaluated at $\Delta\ge\log(2/\eps)$.  W2 is cited here as a \emph{target}, not used.
\end{remark}

\section{The twirled (t-averaged) channels: a mixture}\label{sec:twirl}

\begin{proposition}[The composite of twirled steps is a mixture of corner maps]\label{prop:twirl}
Let $\PP$ be a one-sided protocol satisfying (H) in which the $m$-th step is the \emph{twirled}
zero-collar channel $\beta^{\rm tw}_m=\int\dd\mu(t)\,\beta_{m,t}$, $\beta_{m,t}$ the rotated
zero-collar map with parameter $s_{t}=\lambda_m(1+e^{-2\pi t})$ and $\mu$ the FHSW probability
measure $\dd\mu(t)=\pi(\cosh2\pi t+1)^{-1}\dd t$ (Box~[G5]; $\int\dd\mu=1$ because
$\cosh2\pi t+1=2\cosh^2\pi t$ and $\int\operatorname{sech}^2(\pi t)\dd t=2/\pi$).  Then the recovered state is
the mixture
\begin{equation}\label{eq:twirl}
 \omega\circ\beta^{\rm tw}_1\circ\dots\circ\beta^{\rm tw}_N
 =\int\dd\mu(t_1)\cdots\dd\mu(t_N)\ \omega\circ\beta_{V_{\Phi_{\vec t}}},
\end{equation}
where $\Phi_{\vec t}$ is the composite corner map with corner measure
\begin{equation}\label{eq:twirl-kappa}
 \kappa_k(\vec t)=\sum_{m:\,p^{(m)}=p_k}
 \frac{\lambda_m\bigl(1+e^{-2\pi t_m}\bigr)}{l^{(m)}_B+l^{(m)}_{\rm new}},\qquad
 \zeta^{(I)}_k(\vec t)=\kappa_k(\vec t)\,\bI(p_k),
\end{equation}
and, with $\Phi_{\rm fid}(\vec t):=-\log F(\omega,\omega\circ\beta_{V_{\Phi_{\vec t}}})$,
\begin{equation}\label{eq:twirl-bound}
 \Phi^{\rm tw}_{\rm fid}\ \le\ -\log\int\dd\mu^{\otimes N}(\vec t)\,e^{-\Phi_{\rm fid}(\vec t)}
 \ \le\ \int\dd\mu^{\otimes N}(\vec t)\ \Phi_{\rm fid}(\vec t).
\end{equation}
\end{proposition}
\begin{proof}
\lstep{1}{1}{\eqref{eq:twirl}: composition of maps is bilinear, so the composite of the averages
 is the average of the composites, and $\omega\circ(\cdot)$ is linear.}
\lby{the integrals are $\sigma$-weakly convergent averages of channels with
 $\mu$ a probability measure (Box~[G5]); Fubini for the finite product}
\lstep{1}{2}{\eqref{eq:twirl-kappa}: for fixed $\vec t$ each step is the compression
 \eqref{eq:step-map} with $\sigma_m=s_{t_m}/(l^{(m)}_B+l^{(m)}_{\rm new})$, so
 Theorem~\ref{thm:corner} applies.}
\lby{(C3) and Theorem~\ref{thm:corner}; (H) is a property of the geometry, not of the $s_m$}
\lstep{1}{3}{$F(\omega,\cdot)$ is concave on normal states.}
\lby{Box~[G4]: $P_{\rm tr}(\omega,\psi)=\inf_{a>0}\omega(a)\psi(a^{-1})$ is an infimum of
 affine functions of $\psi$, hence concave; $F=\sqrt{P_{\rm tr}}$ and $\sqrt\cdot$ is concave and
 increasing}
\lstep{1}{4}{\eqref{eq:twirl-bound}: $F\bigl(\omega,\int\omega_{\vec t}\bigr)\ge\int
 F(\omega,\omega_{\vec t})=\int e^{-\Phi_{\rm fid}(\vec t)}$, and
 $-\log\int e^{-\Phi}\dd\mu\le\int\Phi\dd\mu$.}
\lby{$\langle1\rangle3$ (Jensen for the concave $F$), $-\log$ decreasing, and Jensen for the
 convex function $-\log$}
\lqed{1}{5}
\end{proof}

\begin{remark}
We state Proposition~\ref{prop:twirl} and do not develop it: the $\vec t$-integral of
$\Phi_{\rm fid}(\vec t)$ requires the multi-corner law W2 (or, for a single corner,
$\Phi_A$ of paper~1 with $\zeta_t=z(1+e^{-2\pi t})$, Box~[G2]).  Note that
\eqref{eq:twirl-bound} is an upper bound only; the twirled channel could be better than the
average, and at a single corner paper~1's tables already show
$\Phi_A$ convex in $\zeta$, so the mixture bound is not tight.
\end{remark}

\section{Summary}\label{sec:summary}

\begin{table}[h]\centering\small
\caption{Status of every result of this document.  ``PROVED'' means proved here from the cited
results with all hypotheses stated; ``under (H)'', ``under (U)'' name the hypothesis used.}
\label{tab:summary}
\begin{tabular}{@{}lll@{}}
\toprule
result & content & status\\
\midrule
Lem.~\ref{lem:admissible} & $s\ge\lambda\iff k(M)\subseteq A_B$ & PROVED\\
Lem.~\ref{lem:density} & parabolic parameter is a $(-1)$-density, exactly & PROVED\\
Lem.~\ref{lem:field} & first-order field $=\zeta w(y-y_p)\partial_y$, $\zeta=\sigma\bI(p)$;
 Thm.~A's $\zeta=as/(a+L)$ & PROVED (exact)\\
Lem.~\ref{lem:masses}--\ref{lem:jump} & $\Sch=\sum[v]\delta$, no square term; corner mass
 $-2\sigma$ from either side & PROVED\\
Lem.~\ref{lem:compose} & $\Sch(f\circ g)=\Sch(g)+(g')^2\Sch(f)\circ g$ & PROVED\\
Lem.~\ref{lem:same-schwarz} & equal Schwarzians $\iff$ differ by a global M\"obius map & PROVED\\
Thm.~\ref{thm:general} & general transformation rule (masses at $G_m^{-1}(p^{(m)})$) &
 PROVED (no hypothesis)\\
Thm.~\ref{thm:corner} & $\Sch(\Phi)=-2\sum_k\kappa_k\delta_{p_k}$ (W1) & PROVED under (H)\\
Prop.~\ref{prop:H-suff} & sufficient conditions for (H) & PROVED\\
Ex.~\ref{ex:displaced}, \ref{ex:swallow} & (H) fails: mass displaced / corner swallowed;
 W1 wording corrected & PROVED (counterexamples)\\
Lem.~\ref{lem:comp-W}--\ref{lem:moebius-inv} & $V_{k_1}V_{k_2}=V_{k_1k_2}$; $[V_g,P]=0$ for
 $g\in\Mob$ & PROVED\\
Thm.~\ref{thm:state} & composite normal; state depends only on $\Sch(\Phi)$ & PROVED
 (cites Box~[G1])\\
Cor.~\ref{cor:phiA} & a single corner: the state, hence $\Phi_A$, depends only on $\zeta$ &
 PROVED\\
Lem.~\ref{lem:2pt}, Thm.~\ref{thm:rigidity} & vacuum rigidity; no exact recovery (last step keeps a non-empty part) &
 PROVED (no (H) needed)\\
Prop.~\ref{prop:1B1A} & VWZ 1(B) $=$ 1(A) as point maps, every common $t$ & PROVED\\
Prop.~\ref{prop:prot2} & VWZ Protocol 2: corners $p_2,p_3$ & PROVED\\
Thm.~\ref{thm:prot3} & Protocol 3 $=$ over-compressed Theorem-A step, $\zeta_{\rm eff}$,
 $s'>2\lambda'$ & PROVED\\
Cor.~\ref{cor:holonomy} & equal corner measure $\Rightarrow$ equal state & PROVED under (H)\\
Prop.~\ref{prop:startblock}, Lem.~\ref{lem:distinct}, Tab.~\ref{tab:n4}--\ref{tab:n4union} &
 holonomy $=$ starting block; distinct pairs give distinct measures; $n=4,5$ tables & PROVED\\
Lem.~\ref{lem:ampl} & amplification $\zeta^{(I)}=\zeta_{\rm step}(1+z')$, monotone in the frame &
 PROVED\\
Lem.~\ref{lem:LRdata} & L$\to$R chain: $e^{\Delta_k}=(1+u_k)(1+v_k)$;
 $\zeta_k=\frac{\sinh(\Delta_{k+1}/2)}{\sinh(\Delta_k/2)\sinh(\frac{\Delta_k+\Delta_{k+1}}2)}$;
 $e^{\Delta_k}\le1+2/\zeta_k$, equality iff $k=n-1$ & PROVED (new)\\
Thm.~\ref{thm:sep} & all-weak separation $e^{\Delta_k}\ge2/\eps$ (W3) & PROVED; constant sharp
 as $\eps\to0$\\
Cor.~\ref{cor:sep-admissible} & admissible $s_k\ge\lambda_k$: $e^{\Delta_k}\ge1/\eps$ & PROVED\\
Thm.~\ref{thm:w6} & type-III sequential Bures \emph{angle} bound (W6) & PROVED under (U)\\
Thm.~\ref{thm:w6chord} & the same bound for the Bures \emph{distance} & PROVED, no (U)
 (Box~[G7])\\
Cor.~\ref{cor:w6-second} & $\Phi_{\rm tot}\le(\sum\sqrt{\Phi_k})^2(1+o(1))$ & PROVED, no (U);
 the exact form \eqref{eq:w6-master} needs (U)\\
Prop.~\ref{prop:twirl} & twirled composite is a mixture; $\Phi^{\rm tw}\le$ average & PROVED,
 not developed\\
W2 (network law) & second-order kernel $\hat G$ & TARGET (G1b), cited only\\
\bottomrule
\end{tabular}
\end{table}

\paragraph{Numerical support.}
\texttt{rigor/g1a\_calculus.py} ($33$ checks, all PASS; output
\texttt{g1a\_calculus.out}): parabolic group laws; the displaced and the swallowed corner; the
1(B) $=$ 1(A) identity in exact rational arithmetic and for rotated $t$; Protocol 3 in both orders
and its M\"obius reduction; $s'>2\lambda'$ in $20\,000$ rational configurations; the $n=4,5$
protocol tables with a direct finite-difference measurement of every Schwarzian mass; the
amplification identity in $20\,000$ exact rational cases.
\texttt{rigor/g1a\_separation.py} ($15$ checks, all PASS; output
\texttt{g1a\_separation.out}, a verbatim capture of the run): the L$\to$R identities in
$50\,000$ random chains, including \eqref{eq:LR3} in the repaired form ($\le$ always, the exact
gap $2\sinh^2(\Delta_k/2)(\coth(\Delta_{k+1}/2)-1)>0$ for $k\le n-2$, equality at $k=n-1$ in
$100\%$ of the chains); the separation bound in $400\,000$ random chains (minimum
$(\eps/2)e^{\Delta_k}=0.999999996$) and in a constrained minimisation reproducing REF-P6-0's
$20.05$ and $200.01$; the recursion constants; the admissible case; and REF-P6-3's exact
constrained minimum by backward recursion, which reproduces $6.8134$, $20.0499$, $66.6817$,
$200.0050$, $666.6682$ and shows that $6.75$ at $\eps=0.3$ is unattainable.
Re-used without recomputation: \texttt{ref\_p6\_0\_field\_exact.out} (the exact field identity,
$50\,000$ rational cases), \texttt{ref\_p6\_0\_w1\_counterexample.out} (the displaced mass),
\texttt{ref\_p6\_0\_separation.out} (the numerical minima $2/\eps$),
\texttt{ref\_p6\_1\_fig25.out} (why VWZ Fig.~25 measures nothing about $\hat G$; not used here).

\paragraph{What is \emph{not} proved here.}
(i) The second-order network law W2 and anything about the kernel $\hat G$ (track G1b/G3); the
only place where the size of $\Phi_{\rm fid}$ enters is through paper~1's $\Phi_A$ for a
\emph{single} corner.  (ii) Hypothesis~(U) is standard but could not be sourced: Uhlmann 1976 is on disk and does
\emph{not} contain it (Hyp.~\ref{hyp:U}), Araki--Raggio 1982 is NOT OBTAINED.  It is used only
for the \emph{angle} inequality \eqref{eq:w6} and the exact form \eqref{eq:w6-master}; the
second-order bound \eqref{eq:w6-2nd} is unconditional via the chordal route
(Theorem~\ref{thm:w6chord}, Box~[G7]) --- whose own citation, Alberti--Uhlmann Prop.~1(1), is
stated there without proof and attributed to Bures 1969, NOT OBTAINED.  (iii) The separation theorem is
proved for the L$\to$R ordinary-Petz chain (and, with the constant $1/\eps$, for admissible
parameters); we have not attempted the analogous statement for protocols with double corners,
where $\Delta=0$ is allowed.  (iv) Normality of each step is imported from Note~4 Thm.~7.1
(Box~[G1]); this document does not re-prove it.

\section{Repairs after REF-P6-3 (2026-09-16)}\label{sec:repairs}

Referee report \texttt{rigor/referee\_network\_corner\_calculus.tex} (REF-P6-3, 21~pp.; scripts
\texttt{rigor/ref\_p6\_3\_calculus.py}, \texttt{rigor/ref\_p6\_3\_separation.py} with their
\texttt{.out}).  Verdict there: nothing FAILS; two BREAKs, three GAPs and several MINORs.  All of
them are applied here; the locators are of \emph{this} (repaired) version.

\begin{enumerate}[leftmargin=2.2em]
\item \textbf{[BREAK] Thm.~\ref{thm:rigidity}(3)} now carries the hypothesis of (2), ``whose last
 step keeps a non-empty part''; the same qualification was added to the abstract, to the title of
 \S\ref{sec:rigidity} and to Table~\ref{tab:summary}.  As written before, (3) was refuted by
 Remark~\ref{rem:degenerate} (the one-step protocol with $A_B=J_0$, whose $\Phi$ is a global
 M\"obius map and which therefore recovers exactly).
\item \textbf{[BREAK] Box~[G5] and Prop.~\ref{prop:twirl}}: the FHSW weight is
 $p(t)=\pi(\cosh2\pi t+1)^{-1}$ (FHSW eq.~(26), p.~6, now transcribed verbatim with a screenshot),
 not $\frac\pi2(\cdots)$, whose total mass is $\frac12$; with the correct constant $\mu$ is a
 probability measure, \eqref{eq:twirl} is a mixture of states and both Jensen steps of
 Prop.~\ref{prop:twirl} apply.  The box now also records the factor $a$ in FHSW's translation
 eq.~(167), p.~31.
\item \textbf{[GAP] Lem.~\ref{lem:LRdata}, eq.~\eqref{eq:LR3}} now reads
 $e^{\Delta_k}\le1+2/\zeta_k$ with equality iff $k=n-1$ (where $\Delta_n=+\infty$ by
 Def.~\ref{def:LR}), and step $\langle1\rangle10$ splits the two cases; Remark~\ref{rem:sep-meaning}
 was adjusted.  \texttt{g1a\_separation.py} check~(iv), whose multiplicative slack hid the issue,
 was rewritten: it now tests $\le$ everywhere, the \emph{exact} gap
 $2\sinh^2(\Delta_k/2)(\coth(\Delta_{k+1}/2)-1)>0$ for $k\le n-2$, and equality at $k=n-1$.
\item \textbf{[GAP] Lem.~\ref{lem:same-schwarz}} asserts the M\"obius factor has no pole on the
 \emph{open} interval $\Phi_1(J)$ (the closure version is false: $\id$ versus $x/(1-x)$ on
 $(0,1)$), and in Theorem~\ref{thm:state}(2) the clause was deleted as superfluous --- $V_g$ is
 unitary on all of $H$ for every $g\in\Mob$ by Lemma~\ref{lem:moebius-inv}.
\item \textbf{[GAP] Prop.~\ref{prop:startblock}}: the count ``exactly $n-1$ classes'' now has its
 missing step, Lemma~\ref{lem:distinct} (distinct starting pairs give measures that differ at
 $p_{i+1}$ by $\sigma^R_{i+2}>0$); the corner \emph{supports} coincide for all middle pairs and
 therefore cannot separate the classes.
\item \textbf{[MINOR] Hyp.~\ref{hyp:U} and the bibliography}: \texttt{refs/Uhlmann76.pdf} is on
 disk and was read --- the earlier ``both sources NOT OBTAINED'' was false.  It contains no
 standard-form statement, so (U) remains a hypothesis; but Theorem~\ref{thm:w6chord} (new) proves
 the chordal form of W6 from Alberti--Uhlmann Prop.~1(1) (new Box~[G7]) with no (U), and
 Corollary~\ref{cor:w6-second}(2) --- the second-order bound, the form W6 needs --- is now
 unconditional.  Bures~1969 was added to the bibliography as NOT OBTAINED.
\item \textbf{[MINOR] Box~[G4]}: the normalised variational formula is Alberti--Uhlmann
 \emph{Cor.~2(5), p.~8}, not Thm.~2(1) (which is the product form, p.~6); both are now quoted with
 screenshots of the \emph{statements} (the earlier screenshot was p.~11, the proof).
\item \textbf{[MINOR] the verdict box after Thm.~\ref{thm:sep}}: ``sharp'' is now ``sharp as
 $\eps\to0$'', with REF-P6-3's exact constrained minima $6.8134$, $20.0499$, $66.6817$,
 $200.0050$, $666.6682$; REF-P6-0's $\eps=0.3$ entry $6.75$ lies below the true minimum and is
 identified as a float artifact.  \texttt{g1a\_separation.py} now recomputes those minima by the
 backward recursion.
\item \textbf{[MINOR] elsewhere}: the boxed \eqref{eq:corner-calculus} sums over the
 \emph{interior} junctions of $J_N$, and the parenthetical in Theorem~\ref{thm:corner}
 $\langle1\rangle3$ was rewritten; (C3) now states that each step is the \emph{vacuum} Petz map of
 its own inclusion and reconciles $s\in[\lambda,2\lambda]$ ($t\ge0$) with the $t$-average over
 $\mathbb R$ of \S\ref{sec:twirl}; Prop.~\ref{prop:1B1A} notes that 1(A) lives on the coarse chain
 $A\mid B\mid CD$; Theorem~\ref{thm:prot3}(4) now cites Corollary~\ref{cor:phiA} (new) for the
 tacit step ``for a single corner $\Phi_{\rm fid}$ depends only on $\zeta$'', and (C7) says that
 $\Phi_A(\zeta)$ at finite $\zeta$ is a definition; $s'>2\lambda'$ is stated strictly; the verdict
 box in \S\ref{sec:schwarz} says $33$ checks; and \texttt{g1a\_separation.out} is now a verbatim
 capture of the run (with the \texttt{RuntimeWarning} lines it produces).
\end{enumerate}

Not repaired, because they are correct as they stand: the referee's remaining items are
observations about the scripts (the rotated-$t$ checks in \texttt{g1a\_calculus.py} \S4 are
vacuous as independent tests --- the mathematical point, that a \emph{common} $\theta_t$ factors
out, is made in Prop.~\ref{prop:1B1A}$\langle1\rangle3$; and \S5 checks the masses at the claimed
corners but not their absence elsewhere, which under (H) is Theorem~\ref{thm:corner}).  Two
sources remain NOT OBTAINED: Araki--Raggio 1982 and Bures 1969.

\begin{thebibliography}{99}
\bibitem{LX} R.~Longo and F.~Xu, \emph{Relative entropy in CFT}, Adv.\ Math.\ \textbf{337} (2018) 139--170;
 arXiv:1712.07283, \texttt{refs/LongoXu\_1712.07283.pdf}.
\bibitem{BGL} R.~Brunetti, D.~Guido and R.~Longo, \emph{Modular structure and duality in conformal quantum
 field theory}, Commun.\ Math.\ Phys.\ \textbf{156} (1993) 201--219; arXiv:funct-an/9302008,
 \texttt{refs/BGL93\_funct-an-9302008.pdf}.
\bibitem{AU02} P.~M.~Alberti and A.~Uhlmann, \emph{On Bures-distance and $*$-algebraic transition probability
 between inner derived positive linear forms over $W^*$-algebras}, arXiv:math-ph/0202038,
 \texttt{refs/AlbertiUhlmann02\_math-ph-0202038.pdf}.
\bibitem{Uhl76} A.~Uhlmann, \emph{The ``transition probability'' in the state space of a $*$-algebra},
 Rep.\ Math.\ Phys.\ \textbf{9} (1976) 273--279, \texttt{refs/Uhlmann76.pdf} (obtained and read:
 it contains the explicit-formula theorem, p.~277, one-argument concavity of $P_{\rm tr}$
 (eq.~(8), p.~274) and monotonicity under restriction to a subalgebra (eq.~(24), p.~277), but
 \emph{no} standard-form/natural-cone statement and no metric statement).
\bibitem{AR82} H.~Araki and G.~A.~Raggio, \emph{A remark on transition probability}, Lett.\ Math.\ Phys.\
 \textbf{6} (1982) 237--240 (\textbf{NOT OBTAINED}: not in \texttt{refs/}; the statement we would
 need from it is isolated as Hypothesis~\ref{hyp:U}, \S\ref{sec:w6}).
\bibitem{Bures69} D.~Bures, \emph{An extension of Kakutani's theorem on infinite product measures to
 the tensor product of semifinite $W^*$-algebras}, Trans.\ Amer.\ Math.\ Soc.\ \textbf{135} (1969)
 199--212 (\textbf{NOT OBTAINED}: not in \texttt{refs/}; it is the source of the triangle
 inequality for $d_{\rm B}$, which \cite{AU02} Prop.~1(1) states without proof, Box~[G7]).
\bibitem{FHSW} T.~Faulkner, S.~Hollands, B.~Swingle and Y.~Wang, \emph{Approximate recovery and relative
 entropy I}, Commun.\ Math.\ Phys.\ \textbf{389} (2022) 349--397; arXiv:2006.08002,
 \texttt{refs/FHSW20\_2006.08002.pdf}.
\bibitem{VWZ} S.~Vardhan, Z.-W.~Wei and Y.~Zou, \emph{Petz recovery from subsystems in conformal field
 theory}, JHEP \textbf{03} (2024) 016; arXiv:2307.14434, \texttt{refs/VWZ\_2307.14434.pdf}.
\bibitem{PS70} R.~T.~Powers and E.~St\o rmer, \emph{Free states of the canonical anticommutation relations},
 Commun.\ Math.\ Phys.\ \textbf{16} (1970) 1--33, \texttt{refs/PowersStormer70.pdf}.
\bibitem{Choi74} M.-D.~Choi, \emph{A Schwarz inequality for positive linear maps on $C^*$-algebras},
 Illinois J.\ Math.\ \textbf{18} (1974) 565--574 (statement reproved here as
 Lemma~\ref{lem:choi}).
\end{thebibliography}


\subsection*{Residual repairs after REF-P6-3b (orchestrator, 2026-09-16)}
Four wording items from the referee's second pass, applied verbatim: (1) Theorem~\ref{thm:w6chord} and
Corollary~\ref{cor:w6-second} no longer inherit Hypothesis~(U) through ``in the situation of
Theorem~\ref{thm:w6}''; the proof of Theorem~\ref{thm:w6chord} records that only step $\langle1\rangle4$ of
Theorem~\ref{thm:w6} uses (U). (2) The abstract now says ``sharp as $\eps\to0$'' and no longer ``to four
digits''. (3) The abstract advertises the unconditional second-order bound, Corollary~\ref{cor:w6-second}(2).
(4) Proposition~\ref{prop:1B1A}'s range of $\theta_t$ reads $(\tfrac12,1]$ for $t\ge0$; the check script
\texttt{g1a\_separation.py} prints ``attained only as $\eps\to0$''.

\end{document}
